\documentclass[10pt,a4paper]{article}
\usepackage[margin=1.8cm]{geometry}
\usepackage{amsmath,amssymb,enumitem,fontspec,titlesec,multicol}
\setlength{\parindent}{0pt}\setlength{\parskip}{3pt}
\titleformat{\section}{\large\bfseries}{}{0pt}{}
\setlist{nosep,leftmargin=1.6em}
\allowdisplaybreaks
\begin{document}
{\LARGE\bfseries MA1522 Canvas practice quizzes}\par
{\small Semester 1, 2026/27. 37 practice quizzes, 185 questions, as drawn on 24 Nov 2026 (each quiz draws from a pool, so another attempt may show other questions). Questions only: Canvas shows the answers after an attempt is submitted.}\par\medskip

\section{Quiz 1.1}
\textbf{1.} \textit{(choose for each blank)} Consider the following linear system

\[\left\{\begin{array}{rrrrrrr}x_1&+&x_2&-&5x_3&=&2\\ x_1&-&2x_2&+&7x_3&=&3\\ &&-x_2&+&x_3&=&5x_1\\ x_1&&&+&x_3&=&17\\ \end{array}\right.\]

The system is  \underline{\hspace{1.2cm}}\textsubscript{1} 
. It has  \underline{\hspace{1.2cm}}\textsubscript{2} 
 variables and  \underline{\hspace{1.2cm}}\textsubscript{3} 
 equations.

\quad Blank 1: in standard form \textbar{} not in standard form

\quad Blank 2: 1 \textbar{} 2 \textbar{} 3 \textbar{} 4

\quad Blank 3: 1 \textbar{} 2 \textbar{} 3 \textbar{} 4

\smallskip
\textbf{2.} \textit{(choose all that apply)} Which of the following are general solutions to the linear system \(2x-3y+z=1\)?

\begin{enumerate}[label=(\Alph*)]
\item \(x=\frac{1}{2}+\frac{3}{2}s-\frac{1}{2}t,\;y=s, \;z=t\)
\item \(x=s,\; y=-\frac{1}{3}+\frac{2}{3}s+\frac{1}{3}t,\; z=t\)
\item \(x=s,\; y=t,\; z=1-2s+3t\)
\item \(x=\frac{1}{2}+\frac{3}{2}t-\frac{1}{2}s, \;y=s,\;z=t\)
\item \(x=s,\; y=-\frac{1}{3}+\frac{2}{3}t+\frac{1}{3}s,\; z=t\)
\end{enumerate}
\smallskip
\textbf{3.} \textit{(fill in the blanks)} A shop sells three types of products, product X (\$x per unit), product Y (\$y per unit), and product Z (\$z per unit). On a particular day, three people are purchasing these three products from the shop:

\quad\textbullet\ Person A purchases 1 unit of product X, 2 units of product Y, and 3 units of product Z and pays \$140.

\quad\textbullet\ Person B purchases 2 units of product X, 3 units of product Y, and 4 units of product Z and pays \$200.

\quad\textbullet\ Person C purchases 4 units of product X, 1 unit of product Y, and 1 unit of product Z and pays \$90.

Fill in the blanks to construct a linear system representing the shop sales of that day.

\[\left\{\begin{array}{ccccccc} x&+&2y&+&az&=&140\\ bx&+&3y&+&4z&=&200\\ 4x&+&cy&+&z&=&90\end{array}\right.\]

where \(a=\) \underline{\hspace{1.2cm}}\textsubscript{1} 
, \(b=\) \underline{\hspace{1.2cm}}\textsubscript{2} 
, \(c=\) \underline{\hspace{1.2cm}}\textsubscript{3} 
.

\smallskip
\textbf{4.} \textit{(true or false)} Is following linear system

\[\left\{\begin{array}{rrr}2x+y-3&=&0\\x-3y+z&=&0\end{array}\right.\]

homogeneous?

\smallskip
\textbf{5.} \textit{(choose for each blank)} Consider the following linear system

\[\left\{\begin{array}{rrrrrrrrr} 3x_1&+&x_2&-&x_3&&&=&2\\ &&5x_2&+&4x_3&+&x_4&=&-1\\ &&&&-x_3&+&x_4&=&5\\ x_1&&&&&+&2x_4&=&9\\ \end{array}\right.\]

The system is  \underline{\hspace{1.2cm}}\textsubscript{1} 
. It has  \underline{\hspace{1.2cm}}\textsubscript{2} 
 variables and  \underline{\hspace{1.2cm}}\textsubscript{3} 
 equations.

\quad Blank 1: in standard form \textbar{} not in standard form

\quad Blank 2: 1 \textbar{} 2 \textbar{} 3 \textbar{} 4 \textbar{} 5

\quad Blank 3: 1 \textbar{} 2 \textbar{} 3 \textbar{} 4 \textbar{} 5

\smallskip
\section{Quiz 1.2}
\textbf{1.} \textit{(choose one)} Given the following linear system, what is its augmented matrix?

\(\begin{array}{c}x-y+5z=9\\2x+4y=3z\\3x-5=z\end{array}\)

\begin{enumerate}[label=(\Alph*)]
\item \(\left(\begin{array}{ccc|c}1&-1&5&9\\2&4&-3&0\\3&0&-1&5\end{array}\right)\)
\item \(\left(\begin{array}{ccc|c}1&-1&5&9\\2&4&3&0\\3&0&1&5\end{array}\right)\)
\item \(\left(\begin{array}{ccc|c}x&-y&5z&9\\2x&4y&0&3z\\3x&-5&0&z\end{array}\right)\)
\item \(\left(\begin{array}{ccc|c}x&-y&5z&9\\2x&4y&-3z&0\\3x&0y&-z&5\end{array}\right)\)
\end{enumerate}
\smallskip
\textbf{2.} \textit{(choose all that apply)} Which columns of \(\left(\begin{array}{cccccc|c}1&0&-2&0&0&0&1\\0&1&0&0&0&0&5\\0&0&0&0&1&0&2\\0&0&0&0&0&1&-1\end{array}\right)\) are pivot columns?

\begin{enumerate}[label=(\Alph*)]
\item 1
\item 2
\item 3
\item 4
\item 5
\item 6
\item 7
\end{enumerate}
\smallskip
\textbf{3.} \textit{(match)} Match the following augmented matrices to whether it is in row-ecehlon form, reduced row-echelon form, or neither.

\begin{enumerate}[label=(\alph*)]
\item \(\mathbf{A}=\left(\begin{array}{cc|c}1&1&3\\0&3&1\\0&1&0\end{array}\right)\)
\item \(\mathbf{B}=\left(\begin{array}{cc|c}1&0&3\\0&1&-1\\0&0&0\end{array}\right)\)
\item \(\mathbf{C}=\left(\begin{array}{cc|c}1&0&3\\0&0&1\\0&0&0\end{array}\right)\)
\item \(\mathbf{D}=\left(\begin{array}{cccccc|c}1&0&-2&0&0&0&1\\0&1&0&0&0&0&5\\0&0&0&0&1&0&2\\0&0&0&0&0&1&-1\end{array}\right)\)
\end{enumerate}
\quad Options: RREF but not REF \textbar{} Neither \textbar{} REF only \textbar{} RREF

\smallskip
\textbf{4.} \textit{(choose one)} Which of the following augmented matrix corresponds to inconsistent linear system?

\begin{enumerate}[label=(\Alph*)]
\item \(\left(\begin{array}{ccc|c}1131\\03-11\\03-10\end{array}\right)\)
\item \(\left(\begin{array}{ccc|c}1&1&3&1\\0&3&-1&2\\0&0&0&0\end{array}\right)\)
\item \(\left(\begin{array}{ccc|c}1132\\03-11\\0010\end{array}\right)\)
\item \(\left(\begin{array}{ccc|c}1131\\03-12\\0011\end{array}\right)\)
\end{enumerate}
\smallskip
\textbf{5.} \textit{(choose one)} The following is an augmented matrix associated to a linear system in variables \(x,y,z\), \(\left(\begin{array}{ccc|c}0&1&-1&0\\0&0&-1&-2\end{array}\right)\). Which variable should we assigned as a parameter in the general solution?

\begin{enumerate}[label=(\Alph*)]
\item \(x\)
\item \(y\)
\item \(z\)
\end{enumerate}
\smallskip
\section{Quiz 1.3}
\textbf{1.} \textit{(choose one)} Consider the following linear system \[\left\{\begin{array}{rrrrrrrrr} 3x_1&+&x_2&-&x_3&&&=&2\\ &&5x_2&+&4x_3&+&x_4&=&-1\\ &&&&-x_3&+&x_4&=&5\\ x_1&&&&&+&2x_4&=&9\\ \end{array}\right.\] What is the resultant augmented matrix after performing the row operation \(R_2+2R_4\)?

\begin{enumerate}[label=(\Alph*)]
\item \(\left(\begin{array}{cccc|c} 31-102\\ 0541-1\\ 00-115\\ 110847 \end{array}\right)\)
\item \(\left(\begin{array}{cccc|c} 3&1&-1&0&2\\ 1&10&8&4&7\\ 0&0&-1&1&5\\ 0&5&4&1&-1 \end{array}\right)\)
\item \(\left(\begin{array}{cccc|c} 31-102\\ 254517\\ 00-115\\ 10029 \end{array}\right)\)
\item \(\left(\begin{array}{cccc|c} 31-102\\ 10029\\ 00-115\\ 254517 \end{array}\right)\)
\end{enumerate}
\smallskip
\textbf{2.} \textit{(choose one)} Which of the following is an elementary row operation?

\begin{enumerate}[label=(\Alph*)]
\item Add \(a\) times of row 1 to row 2, for any real number \(a\).
\item Multiplying row 2 by the constant 0.
\item Add 2 times of row 1 to 3  times of row 2.
\end{enumerate}
\smallskip
\textbf{3.} \textit{(choose one)} Given the following linear system

\(\left\{\begin{array}{c}x-y+5z=9\\x+3y-z=1\\x-3z=1\end{array}\right.\)

what is the resultant linear system after we add 2 times of row 2 to row 1?

\begin{enumerate}[label=(\Alph*)]
\item \(x-y+5z=9\\ 3x+5y+3z=11\\x-3z=1\)
\item \(3x+y+9z=19\\ x+3y-z=1\\ x-3z=1\)
\item \(x-y+5z=9\\3x+y+9z=19\\x-3z=1\)
\item \(3x+5y+3z=11\\x+3y-z=1\\x-3z=1\)
\end{enumerate}
\smallskip
\textbf{4.} \textit{(choose one)} The \textbf{resultant augmented matrix} after performing the row operation \(R_2+3R_1\) is \[\left(\begin{array}{ccc|c}      3  &   3  &   0  &   3\\      3  &   2  &   1  &  -1\\     -1  &  -1  &   3  &   3 \end{array}\right).\]

Which of the following is the original augmented matrix?

\begin{enumerate}[label=(\Alph*)]
\item \(\left(\begin{array}{ccc|c}      3  &   3  &   0  &   3\\     12  &  11  &   1  &   8\\     -1  &  -1  &   3  &   3 \end{array}\right).\)
\item \(\left(\begin{array}{ccc|c}      3  &   3  &   0  &   3\\      2  &   1  &   1  &  -2\\     -1  &  -1  &   3  &   3 \end{array}\right).\)
\item 
\item 
\end{enumerate}
\smallskip
\textbf{5.} \textit{(choose one)} Which of the following is the ‘reverse’ of the elementary row operation where <span class="cke\_widget\_wrapper cke\_widget\_inline cke\_widget\_katex cke\_widget\_wrapper\_math-tex cke\_widget\_selected" role="region" data-cke-widget-wrapper="1" data-cke-filter="off" data-cke-display-name="math" data-cke-widget-id="7" data-math="\(2\)" data-math-display-mode="false" aria-label="math widget">2 times of row <span class="cke\_widget\_wrapper cke\_widget\_inline cke\_widget\_katex cke\_widget\_wrapper\_math-tex cke\_widget\_selected" role="region" data-cke-widget-wrapper="1" data-cke-filter="off" data-cke-display-name="math" data-cke-widget-id="6" data-math="\(1\)" data-math-display-mode="false" aria-label="math widget">1 is added to row <span class="cke\_widget\_wrapper cke\_widget\_inline cke\_widget\_katex cke\_widget\_wrapper\_math-tex cke\_widget\_selected" role="region" data-cke-widget-wrapper="1" data-cke-filter="off" data-cke-display-name="math" data-cke-widget-id="5" data-math="\(2\)" data-math-display-mode="false" aria-label="math widget">2, \(R_2+2R_1\)?

\begin{enumerate}[label=(\Alph*)]
\item Add 2 times of row 1 to row 2, \(R_2+2R_1\)
\item Add -2 times of row 1 to row 2, \(R_2-2R_1\)
\item Add \(1/2\) times of row 1 to row 2, \(R_2+\frac{1}{2}R_1\)
\end{enumerate}
\smallskip
\section{Quiz 1.4}
\textbf{1.} \textit{(choose one)} In performing the Gaussian Elimination, what should be the next step for the following augmented matrix \[\left(\begin{array}{ccc|c}      0&0&0&0\\      0&0&1&-1\\      0&2&3&-2\\      0&1&-1&2 \end{array}\right)?\]

\begin{enumerate}[label=(\Alph*)]
\item \(R_1\leftrightarrow R_4\)
\item \(R_2\leftrightarrow R_4\)
\item \(R_1\leftrightarrow R_2\)
\item \(R_3\leftrightarrow R_4\)
\end{enumerate}
\smallskip
\textbf{2.} \textit{(choose one)} Consider the following augmented matrix

\[\left(\begin{array}{ccc|c} 1&2&3&26\\ 4&5&6&62\\ 7&8&9&98 \end{array}\right)\]

Which of the following matrices in row-echelon form is \textbf{not} row equivalent to the augmented matrix above?

\begin{enumerate}[label=(\Alph*)]
\item \[\left(\begin{array}{ccc|c} 12011\\03642\\0015\end{array}\right)\]
\item \[\left(\begin{array}{ccc|c}12326\\01214\\0000\end{array}\right)\]
\item \[\left(\begin{array}{ccc|c}45662\\03642\\0000\end{array}\right)\]
\item \[\left(\begin{array}{ccc|c}10-1-2\\01214\\0000\end{array}\right)\]
\end{enumerate}
\smallskip
\textbf{3.} \textit{(choose one)} In performing the Gauss-Jordan Elimination, what should be the next step for the following augmented matrix \[\left(\begin{array}{ccc|c}      2&1&1&-1\\      0&1&3&-2\\      0&0&1&-1\\      0&0&0&0 \end{array}\right)?\]

\begin{enumerate}[label=(\Alph*)]
\item \(\frac{1}{2}R_1\)
\item \(R_1+R_2\)
\item \(R_2-\frac{1}{2}R_1\)
\item \(R_3\leftrightarrow R_2\)
\end{enumerate}
\smallskip
\textbf{4.} \textit{(choose one)} In performing the Gaussian Elimination, what should be the next step for the following augmented matrix

\[\left(\begin{array}{ccc|c}     0&2&1&-1\\     0&1&3&-2\\     0&1&0&-1\\     0&1&-1&2\end{array}\right)?\]

\begin{enumerate}[label=(\Alph*)]
\item \(R_2-\frac{1}{2}R_1\)
\item \(R_2- 2R_1\)
\item \(R_1\leftrightarrow R_2\)
\item \(R_3\leftrightarrow R_1\)
\end{enumerate}
\smallskip
\textbf{5.} \textit{(choose one)} In performing the Gaussian Elimination, what should be the next step for the following augmented matrix \[\left(\begin{array}{ccc|c}      0&0&1&-1\\      0&2&3&-2\\      0&1&0&-1\\      0&1&-1&2 \end{array}\right)?\]

\begin{enumerate}[label=(\Alph*)]
\item \(R_1\leftrightarrow R_2\)
\item \(R_1+R_2\)
\item \(R_2-R_1\)
\item \(R_3\leftrightarrow R_2\)
\end{enumerate}
\smallskip
\section{Quiz 1.5}
\textbf{1.} \textit{(fill in the blanks)} Consider the following linear system

\(\left\{\begin{array}{ccccc}&x&+&y&+&z&=&1\\ &&&ay&-&z&=&0\\&&&&& (a-1)z&=&a \end{array}\right.\)

For the system to have no solution, \(a=\) \underline{\hspace{1.2cm}}\textsubscript{1} 
.

For the system to have infinitely many solutions, \(a=\) \underline{\hspace{1.2cm}}\textsubscript{2} 
.

\smallskip
\textbf{2.} \textit{(choose one)} What are the cases to consider for the unknown \(a\) in solving the following augmented matrix? \[\left(\begin{array}{cccc|c} 1&a+1&2&-1&1\\ 0&a-2&1&2&-3\\ 0&0&1&a-3&1\\ 0&0&0&1&a+4 \end{array}\right)?\]

\begin{enumerate}[label=(\Alph*)]
\item 2
\item 1
\item 3
\item 4
\item None of the other options above.
\end{enumerate}
\smallskip
\textbf{3.} \textit{(choose one)} Which step should we perform to reduce the following augmented matrix? \[\left(\begin{array}{ccc|c} 1&0&2&-1\\ a&0&1&2\\ 0&1&3&5\\ a&1&-2&1 \end{array}\right)?\]

\begin{enumerate}[label=(\Alph*)]
\item \(R_2-aR_1\)
\item \(\frac{1}{a}R_2\)
\item \(R_2\leftrightarrow R_1\)
\item None of the other options. We should be considering the cases for \(a\) instead.
\end{enumerate}
\smallskip
\textbf{4.} \textit{(choose one)} Consider the following linear system, where \(a\) is a real constant.

\[\left\{\begin{array}{ccccccc} x&+&2y&+&2z&=&6\\ 2x&+&3y&+&4z&=&15\\ 3x&+&5y&+&6z&=&a \end{array}\right.\]

For what values of \(a\) will the system has infinitely many solutions?

\begin{enumerate}[label=(\Alph*)]
\item 21
\item 14
\item 7
\item 28
\end{enumerate}
\smallskip
\textbf{5.} \textit{(choose all that apply)} Which of the following statements is/are true about the reduce row-echelon form of a matrix?

\begin{enumerate}[label=(\Alph*)]
\item The number of pivot columns is equal to the number of nonzero rows.
\item The number of non-pivot columns is equal to the number of zero rows.
\item The number of pivot columns is equal to the number of zero rows.
\item The number of non-pivot columns is equal to the number of nonzero rows.
\end{enumerate}
\smallskip
\section{Quiz 2.1}
\textbf{1.} \textit{(choose one)} For \(\begin{pmatrix}1&-1&0\\-1&1&2\\a&2&3\end{pmatrix}\) to be symmetric, \(a\) must be

\begin{enumerate}[label=(\Alph*)]
\item 0
\item -1
\item 1
\item 2
\end{enumerate}
\smallskip
\textbf{2.} \textit{(choose one)} Let \(\mathbf{A}\) be a \(n\times n\) diagonal matrix. Which of the following row operations performed on \(\mathbf{A}\) will still result in a diagonal matrix?

\begin{enumerate}[label=(\Alph*)]
\item \(aR_i\) for some real number \(a\) and \(i=1,...,n\).
\item \(R_i\leftrightarrow R_j\) for some \(i,j=1,...,n\).
\item \(R_i+cR_j\) for some \(i,j=1,...,n\).
\end{enumerate}
\smallskip
\textbf{3.} \textit{(choose one)} Which of the following statements is true?

\begin{enumerate}[label=(\Alph*)]
\item A scalar matrix is a diagonal matrix.
\item A diagonal matrix is a scalar matrix.
\end{enumerate}
\smallskip
\textbf{4.} \textit{(choose all that apply)} The matrix

\[\mathbf{A}=\begin{pmatrix}0&1&-1&3\\0&0&2&17\\0&0&0&-3\\0&0&0&0\end{pmatrix}\]

is a (an)

\begin{enumerate}[label=(\Alph*)]
\item upper triangular matrix
\item strictly upper triangular matrix
\item lower triangular matrix
\item strictly lower triangular matrix
\item None of the other options are right.
\end{enumerate}
\smallskip
\textbf{5.} \textit{(choose all that apply)} A zero square matrix is a

\begin{enumerate}[label=(\Alph*)]
\item diagonal matrix
\item scalar matrix
\item symmetric matrix
\item triangular matrix
\end{enumerate}
\smallskip
\section{Quiz 2.2}
\textbf{1.} \textit{(choose one)} Which of the following statements is true about symmetric matrices?

\begin{enumerate}[label=(\Alph*)]
\item If \(\mathbf{A}\) is symmetric, then \(\mathbf{A}^2\) is symmetric too.
\item If \(\mathbf{A}^2\) is symmetric, then necessarily \(\mathbf{A}\) is symmetric.
\end{enumerate}
\smallskip
\textbf{2.} \textit{(choose one)} Which of the following is the transpose of post-multiplying \(\begin{pmatrix}1&5&3\\0&0&1\\1&2&0\end{pmatrix}\) to matrix \(\begin{pmatrix}2&1&5\\1&-2&1\\1&-1&0\end{pmatrix}\)?

\begin{enumerate}[label=(\Alph*)]
\item \(\begin{pmatrix}7&2&1\\20&7&5\\7&1&2\end{pmatrix}\)
\item \(\begin{pmatrix}10&1&4\\-12&-1&3\\10&0&7\end{pmatrix}\)
\item \(\begin{pmatrix}7&20&7\\2&7&1\\1&5&2\end{pmatrix}\)
\item \(\begin{pmatrix}10&-12&10\\1&-1&0\\4&3&7\end{pmatrix}\)
\end{enumerate}
\smallskip
\textbf{3.} \textit{(choose all that apply)} Let \(\mathbf{A}\) be a \(m\times n\) matrix. Which of the following matrix multiplication is well-defined?

\begin{enumerate}[label=(\Alph*)]
\item \(\mathbf{AA}^T\)
\item \(\mathbf{A}^T\mathbf{A}\)
\item \(\mathbf{AA}\)
\item \(\mathbf{A}^T\mathbf{A}^T\)
\end{enumerate}
\smallskip
\textbf{4.} \textit{(choose one)} Let \(\mathbf{A}=\begin{pmatrix}1&0&1\\1&1&0\\0&1&1\end{pmatrix}\). Which of the following is \(\mathbf{A}^2\)?

\begin{enumerate}[label=(\Alph*)]
\item \(\begin{pmatrix}1&1&2\\2&1&1\\1&2&1\end{pmatrix}\)
\item \(\frac{1}{4}\begin{pmatrix}-1&3&-1\\-1&-1&3\\3&-1&-1\end{pmatrix}\)
\item \(\frac{1}{2}\begin{pmatrix}1&1&-1\\-1&1&1\\1&-1&1\end{pmatrix}\)
\item \(\begin{pmatrix}1&0&1\\1&1&0\\0&1&1\end{pmatrix}\)
\end{enumerate}
\smallskip
\textbf{5.} \textit{(fill in the blanks)} Evaluate the following.

\(3\begin{pmatrix}1\\-1\\3\end{pmatrix}-2\begin{pmatrix}1\\1\\2\end{pmatrix}=\begin{pmatrix}a\\b\\c\end{pmatrix}\), where

\(a=\) \underline{\hspace{1.2cm}}\textsubscript{1} 
, \(b=\) \underline{\hspace{1.2cm}}\textsubscript{2} 
, \(c=\) \underline{\hspace{1.2cm}}\textsubscript{3}

\smallskip
\section{Quiz 2.3}
\textbf{1.} \textit{(fill in the blanks)} Suppose a linear system, expressed as a matrix equation, is

\(\begin{pmatrix}1&a&2\\3&-1&2\\2&3&b\\3&3&0\end{pmatrix}\begin{pmatrix}x\\y\\z\end{pmatrix}=\begin{pmatrix}2\\1\\3\\1\end{pmatrix},\)

and expressed as a vector equation, is

\(x\begin{pmatrix}1\\3\\2\\3\end{pmatrix}+y\begin{pmatrix}2\\-1\\c\\3\end{pmatrix}+z\begin{pmatrix}2\\2\\2\\d\end{pmatrix}=\begin{pmatrix}2\\1\\3\\1\end{pmatrix}\).

Then \(a=\) \underline{\hspace{1.2cm}}\textsubscript{1} 
, \(b=\) \underline{\hspace{1.2cm}}\textsubscript{2} 
, \(c=\) \underline{\hspace{1.2cm}}\textsubscript{3} 
, and \(d=\) \underline{\hspace{1.2cm}}\textsubscript{4} 
.

\smallskip
\textbf{2.} \textit{(choose one)} Which of the following equality is true?

\begin{enumerate}[label=(\Alph*)]
\item \(\mathbf{A}\begin{pmatrix}\mathbf{b}_1&\mathbf{b}_2&\cdots&\mathbf{b}_n\end{pmatrix}=\begin{pmatrix}\mathbf{Ab}_1&\mathbf{Ab}_2&\cdots&\mathbf{Ab}_n\end{pmatrix}\)
\item \(\begin{pmatrix}\mathbf{b}_1&\mathbf{b}_2&\cdots&\mathbf{b}_n\end{pmatrix}\mathbf{A}=\begin{pmatrix}\mathbf{b}_1\mathbf{A}&\mathbf{b}_2\mathbf{A}&\cdots&\mathbf{b}_n\mathbf{A}\end{pmatrix}\)
\item \(\mathbf{B}\begin{pmatrix}\mathbf{a}_1\\\mathbf{a}_2\\\vdots\\\mathbf{a}_m\end{pmatrix}=\begin{pmatrix}\mathbf{Ba}_1\\\mathbf{Ba}_2\\\vdots\\\mathbf{Ba}_m\end{pmatrix}\)
\end{enumerate}
\smallskip
\textbf{3.} \textit{(fill in the blanks)} If \(\mathbf{A}=\begin{pmatrix}1&-3\\-3&8\end{pmatrix}\) and \(\mathbf{AB}=\begin{pmatrix}-1&3&-2\\1&-7&3\end{pmatrix}\).

The second column of \(\mathbf{B}\) is \(\begin{pmatrix}a\\b\end{pmatrix}\), where \(a=\) \underline{\hspace{1.2cm}}\textsubscript{1} 
, and \(b=\) \underline{\hspace{1.2cm}}\textsubscript{2}

\smallskip
\textbf{4.} \textit{(choose one)} Let \(\mathbf{A}\) be a \(m\times p\) matrix and \(\mathbf{B}\) a \(p\times n\) matrix. Write \(\mathbf{B}\) in terms of its columns, \(\mathbf{B}=\begin{pmatrix}\mathbf{b}_1&\mathbf{b}_2&\cdots&\mathbf{b}_n\end{pmatrix}\). Write \(\mathbf{x}=(x_i)\) as a column vector. Then

\begin{enumerate}[label=(\Alph*)]
\item \(\mathbf{ABx}=x_1\mathbf{Ab}_1+x_2\mathbf{Ab}_2+\cdots+x_n\mathbf{Ab}_n\)
\item \(\mathbf{ABx}=\mathbf{Ab}_1\mathbf{x}+\mathbf{Ab}_2\mathbf{x}+\cdots+\mathbf{Ab}_n\mathbf{x}\)
\item \(\mathbf{ABx}=\mathbf{b}_1\mathbf{Ax}+\mathbf{b}_2\mathbf{Ax}+\cdots+\mathbf{b}_n\mathbf{Ax}\)
\item \(\mathbf{ABx}=x_1\mathbf{b}_1\mathbf{A}+x_2\mathbf{b}_2\mathbf{A}+\cdots+x_n\mathbf{b}_n\mathbf{A}\)
\end{enumerate}
\smallskip
\textbf{5.} \textit{(choose one)} Which equation represents \(\left\{\begin{array}{c}2x-y=3\\x+y=-1\end{array}\right.\)?

\begin{enumerate}[label=(\Alph*)]
\item \(x\begin{pmatrix}2\\1\end{pmatrix}+y\begin{pmatrix}-1\\1\end{pmatrix}=\begin{pmatrix}3\\-1\end{pmatrix}\)
\item \(\begin{pmatrix}2&-1\\1&1\end{pmatrix}\begin{pmatrix}3\\-1\end{pmatrix}=\begin{pmatrix}x\\y\end{pmatrix}\)
\item \(x\begin{pmatrix}2\\-1\end{pmatrix}+y\begin{pmatrix}1\\1\end{pmatrix}=\begin{pmatrix}3\\-1\end{pmatrix}\)
\item \(3\begin{pmatrix}2\\1\end{pmatrix}-\begin{pmatrix}-1\\1\end{pmatrix}=\begin{pmatrix}x\\y\end{pmatrix}\)
\end{enumerate}
\smallskip
\section{Quiz 2.4}
\textbf{1.} \textit{(choose one)} Let \(\mathbf{A}=\begin{pmatrix}1&0&1&0\\1&1&0&0\\0&1&1&0\end{pmatrix}\) and \(\mathbf{B}=\begin{pmatrix}1/2&1/2&-1/2\\-1/2&1/2&1/2\\1/2&-1/2&1/2\\0&0&0\end{pmatrix}\). Given that \(\mathbf{AB}=\begin{pmatrix}1&0&0\\0&1&0\\0&0&1\end{pmatrix}\). Which of the following statements is true?

\begin{enumerate}[label=(\Alph*)]
\item \(\mathbf{A}\) is not invertible.
\item \(\mathbf{A}\) is the inverse of \(\mathbf{B}\).
\item \(\mathbf{A}\) is the only matrix such that \(\mathbf{AB=I}\).
\item \(\mathbf{BA=I}\) too.
\end{enumerate}
\smallskip
\textbf{2.} \textit{(choose one)} Suppose \(\begin{pmatrix}1\\-1\end{pmatrix}\) is a solution to \(\mathbf{Ax=0}\) for an 2 by 2 square matrix \(\mathbf{A}\). Which statement is false?

\begin{enumerate}[label=(\Alph*)]
\item If \(\mathbf{Ax=}\begin{pmatrix}a\\b\end{pmatrix}\) is consistent, then the solution may be unique.
\item \(\mathbf{A}\) is not invertible.
\item For any vector \(\mathbf{B}\), either \(\mathbf{Ax=b}\) is inconsistent, or it has infinitely many solutions.
\item \(\mathbf{Ax=0}\) have infinitely many solutions.
\end{enumerate}
\smallskip
\textbf{3.} \textit{(choose one)} Let \(\mathbf{A}\) and \(\mathbf{B}\)  be square matrices of order \(n\). If \(\mathbf{A}\mathbf{B}\) is invertible, then both \(\mathbf{A}\) and \(\mathbf{B}\) must be invertible.

\begin{enumerate}[label=(\Alph*)]
\item True
\item False
\end{enumerate}
\smallskip
\textbf{4.} \textit{(choose one)} The inverse of an invertible symmetric matrix is symmetric.

\begin{enumerate}[label=(\Alph*)]
\item True
\item False
\end{enumerate}
\smallskip
\textbf{5.} \textit{(choose one)} Which of the following is the inverse of the matrix \(\begin{pmatrix}1&0&1\\1&1&0\\0&1&1\end{pmatrix}\)?

\begin{enumerate}[label=(\Alph*)]
\item \(\frac{1}{2}\begin{pmatrix}1&1&-1\\-1&1&1\\1&-1&1\end{pmatrix}\)
\item \(\begin{pmatrix}1&1&0\\0&1&1\\1&0&1\end{pmatrix}\)
\item \(\begin{pmatrix}1&1&-1\\-1&1&1\\1&-1&1\end{pmatrix}\)
\item \(\frac{1}{2}\begin{pmatrix}1&1&0\\0&1&1\\1&0&1\end{pmatrix}\)
\end{enumerate}
\smallskip
\section{Quiz 2.5}
\textbf{1.} \textit{(match)} Match the elementary row operations to the order 4 elementary matrices.

Elementary matrix 1: \(\begin{pmatrix}1&0&0&1\\0&1&0&0\\0&0&1&0\\0&0&0&1\end{pmatrix}\)

Elementary matrix 2: \(\begin{pmatrix}1&0&0&0\\0&1&0&0\\0&0&4&0\\0&0&0&1\end{pmatrix}\)

Elementary matrix 3: \(\begin{pmatrix}0&0&0&1\\0&1&0&0\\0&0&1&0\\1&0&0&0\end{pmatrix}\)

Elementary matrix 4: \(\begin{pmatrix}1&0&0&0\\0&1&0&0\\0&0&1&0\\0&-1&0&1\end{pmatrix}\)

Elementary matrix 5: \(\begin{pmatrix}1&0&0&0\\0&1&0&0\\0&0&1&0\\0&0&0&3\end{pmatrix}\)

Elementary matrix 6: \(\begin{pmatrix}1&0&0&0\\0&1&0&-c\\0&0&1&0\\0&0&0&1\end{pmatrix}\)

Elementary matrix 7: \(\begin{pmatrix}1&0&0&0\\0&1&0&0\\0&0&1&0\\1&0&0&1\end{pmatrix}\)

\begin{enumerate}[label=(\alph*)]
\item \(R_2-cR_4\)
\item \(R_1\leftrightarrow R_4\)
\item \(4R_3\)
\item \(R_1+R_4\)
\end{enumerate}
\quad Options: Elementary matrix 5 \textbar{} Elementary matrix 6 \textbar{} Elementary matrix 1 \textbar{} Elementary matrix 7 \textbar{} Elementary matrix 3 \textbar{} Elementary matrix 2 \textbar{} Elementary matrix 4

\smallskip
\textbf{2.} \textit{(choose one)} Suppose \(\mathbf{A}\) is a \(3\times 4\) matrix. Then the elementary matrix corresponding to the row operations \(R_3+2R_2\) is

\begin{enumerate}[label=(\Alph*)]
\item \(\begin{pmatrix}1&0&0\\0&1&0\\0&2&1\end{pmatrix}\)
\item \(\begin{pmatrix}1&0&0\\0&1&2\\0&0&1\end{pmatrix}\)
\item \(\begin{pmatrix}1&0&0&0\\0&1&0&0\\0&2&1&0\\0&0&0&1\end{pmatrix}\)
\item \(\begin{pmatrix}1&0&0&0\\0&1&2&0\\0&0&1&0\\0&0&0&1\end{pmatrix}\)
\end{enumerate}
\smallskip
\textbf{3.} \textit{(choose one)} Suppose \[\left( \begin{array}{ccc} 1&0&0\\0&1&0\\0&0&1\end{array} \right)\xrightarrow{R_3- R_1}\xrightarrow{R_2-2R_1}\xrightarrow{R_2\leftrightarrow R_3}\mathbf{A}.\] Write \(\mathbf{A}\) as a product of elementary row operations.

\begin{enumerate}[label=(\Alph*)]
\item \(\mathbf{A}=\begin{pmatrix}1&0&0\\0&0&1\\0&1&0\end{pmatrix}\begin{pmatrix}1&0&0\\-2&1&0\\0&0&1\end{pmatrix}\begin{pmatrix}1&0&0\\0&1&0\\-1&0&1\end{pmatrix}\)
\item \(\mathbf{A}=\begin{pmatrix}1&0&0\\0&0&1\\0&1&0\end{pmatrix}\begin{pmatrix}1&0&0\\0&1&0\\-1&0&1\end{pmatrix}\begin{pmatrix}1&0&0\\-2&1&0\\0&0&1\end{pmatrix}\)
\item \(\mathbf{A}=\begin{pmatrix}1&0&0\\0&0&1\\0&1&0\end{pmatrix}\begin{pmatrix}1&-2&0\\0&1&0\\0&0&1\end{pmatrix}\begin{pmatrix}1&0&-1\\0&1&0\\0&0&1\end{pmatrix}\)
\item \(\mathbf{A}=\begin{pmatrix}1&0&0\\0&0&1\\0&1&0\end{pmatrix}\begin{pmatrix}1&0&-1\\0&1&0\\0&0&1\end{pmatrix}\begin{pmatrix}1&-2&0\\0&1&0\\0&0&1\end{pmatrix}\)
\item None of the other options are right.
\end{enumerate}
\smallskip
\textbf{4.} \textit{(match)} Match the following elementary row operations to its inverses.

Elementary matrix 1: \(\begin{pmatrix}1&0&0\\2&1&0\\0&0&1\end{pmatrix}\)

Elementary matrix 2: \(\begin{pmatrix}1&2&0\\0&1&0\\0&0&1\end{pmatrix}\)

Elementary matrix 3: \(\begin{pmatrix}1&0&0\\0&1&0\\0&0&-3\end{pmatrix}\)

Elementary matrix 4: \(\begin{pmatrix}1&1/2&0\\0&1&0\\0&0&1\end{pmatrix}\)

Elementary matrix 5: \(\begin{pmatrix}1&0&0\\0&1&0\\0&0&1/3\end{pmatrix}\)

Elementary matrix 6: \(\begin{pmatrix}0&1&0\\1&0&0\\0&0&1\end{pmatrix}\)

\begin{enumerate}[label=(\alph*)]
\item \(\begin{pmatrix}0&1&0\\1&0&0\\0&0&1\end{pmatrix}\)
\item \(\begin{pmatrix}1&0&0\\0&1&0\\0&0&3\end{pmatrix}\)
\item \(\begin{pmatrix}1&-2&0\\0&1&0\\0&0&1\end{pmatrix}\)
\end{enumerate}
\quad Options: Elementary matrix 6 \textbar{} Elementary matrix 3 \textbar{} Elementary matrix 4 \textbar{} Elementary matrix 1 \textbar{} Elementary matrix 2 \textbar{} Elementary matrix 5

\smallskip
\textbf{5.} \textit{(choose one)} What is \(\begin{pmatrix}1&0&3\\0&1&0\\0&0&1\end{pmatrix}^{-1}\)?

\begin{enumerate}[label=(\Alph*)]
\item \(\begin{pmatrix}1&0&-3\\0&1&0\\0&0&1\end{pmatrix}\)
\item \(\begin{pmatrix}1&0&\frac{1}{3}\\0&1&0\\0&0&1\end{pmatrix}\)
\item \(\begin{pmatrix}1&0&-\frac{1}{3}\\0&1&0\\0&0&1\end{pmatrix}\)
\item \(\begin{pmatrix}1&0&0\\0&1&0\\3&0&1\end{pmatrix}\)
\item \(\begin{pmatrix}1&0&0\\0&1&0\\-3&0&1\end{pmatrix}\)
\end{enumerate}
\smallskip
\section{Quiz 2.6}
\textbf{1.} \textit{(choose one)} Suppose \(\mathbf{A}\) is a 3 by 3 matrix such that \(\mathbf{BA=I}\) for some matrix \(\mathbf{B}\). Then the system \(\mathbf{Ax=b}\)

\begin{enumerate}[label=(\Alph*)]
\item has only a unique solution
\item has infinitely many solutionv
\item may be inconsistent
\end{enumerate}
\smallskip
\textbf{2.} \textit{(choose one)} Suppose \(\mathbf{A}\) is a \(n\times m\) matrix such that there exists a \(\mathbf{B}\) with \(\mathbf{AB=BA}\). What can we conclude?

\begin{enumerate}[label=(\Alph*)]
\item \(\mathbf{A}\) must be a square matrix.
\item \(\mathbf{A}\) is invertible.
\item \(\mathbf{B}\) is invertible.
\end{enumerate}
\smallskip
\textbf{3.} \textit{(choose one)} Suppose \(\mathbf{A}\) is a square matrix such that there exists some matrix \(\mathbf{B}\) with \(\mathbf{AB=I}\). Which statement is false?

\begin{enumerate}[label=(\Alph*)]
\item The matrix \(\mathbf{B}\) is not necessarily unique.
\item It must be that \(\mathbf{BA=I}\).
\item Any row-echelon form of \(\mathbf{A}\) do not have non-pivot columns
\item The reduced row-echelon form of \(\mathbf{A}\) is the identity matrix.
\end{enumerate}
\smallskip
\textbf{4.} \textit{(choose one)} Let \(\mathbf{A}\) be a 5 by 5 matrix and suppose \(\mathbf{Ax}=\begin{pmatrix}2\\1\\5\\-1\\3\end{pmatrix}\) is inconsistent. Then the homogeneous system \(\mathbf{Ax=0}\) must have infinitely many solution.

\begin{enumerate}[label=(\Alph*)]
\item True
\item False
\end{enumerate}
\smallskip
\textbf{5.} \textit{(choose one)} Suppose \(\mathbf{A}\) is a \(3\times 3\) invertible matrix and \(\mathbf{b}\) a \(3\times 1\) vector. Consider the linear system \(\mathbf{Ax=b}\). Which of the following statements is true?

\begin{enumerate}[label=(\Alph*)]
\item If \(\mathbf{b=0}\), then the 3 planes representing the 3 equations must intersect at the origin.
\item If \(\mathbf{b\neq 0}\), then the 3 planes representing the 3 equations may not intersect.
\item If \(\mathbf{b=0}\), then the 3 planes representing the 3 equations may intersection along a line.
\end{enumerate}
\smallskip
\section{Quiz 2.7}
\textbf{1.} \textit{(choose one)} Suppose

\[\mathbf{A}\xrightarrow{R_2+R_1}\xrightarrow{R_3+R_1}\xrightarrow{R_3-2R_2}\begin{pmatrix}     -1  &  -1  &   0  &   0\\      0  &  -1  &   0  &   0\\      0  &   0  &   1  &   2 \end{pmatrix}.\]

Write \(\mathbf{A=LU}\), where \(\mathbf{U}=\begin{pmatrix}     -1  &  -1  &   0  &   0\\      0  &  -1  &   0  &   0\\      0  &   0  &   1  &   2 \end{pmatrix}\).

What is the \((3,1)\) entry of \(\mathbf{L}\)?

\begin{enumerate}[label=(\Alph*)]
\item -1
\item 1
\item 2
\item We do not have enough information to know the \((3,1)\) entry of \(\mathbf{L}\).
\item -2
\item 0
\end{enumerate}
\smallskip
\textbf{2.} \textit{(choose one)} Let \(\mathbf{A=LU}\) be a LU-factorization of \(\mathbf{A}\). Which of the following statements is true?

\begin{enumerate}[label=(\Alph*)]
\item \(\mathbf{A}\) is invertible if and only if \(\mathbf{U}\) is invertible.
\item \(\mathbf{A}\) is invertible if and only if \(\mathbf{L}\) is invertible.
\item The invertibility of \(\mathbf{A}\) is not dependent on \(\mathbf{L}\) nor \(\mathbf{U}\).
\end{enumerate}
\smallskip
\textbf{3.} \textit{(choose one)} It is given that \(\mathbf{A}=\begin{pmatrix}      3   &  0   &  3   &  0\\      3   &  1   &  3   &  1\\      2   & -1  &   3  &  -1\\ \end{pmatrix}\) has a LU factorization \(\mathbf{A=LU}\), where

\[\mathbf{U}=\begin{pmatrix}      3   &  0   &  3  &   0\\      0   &  1   &  0  &   1\\      0   &  0   &  1  &   0 \end{pmatrix}.\]

What is the \((3,1)\)-entryof \(\mathbf{L}\)?

\begin{enumerate}[label=(\Alph*)]
\item 2/3
\item 1/3
\item -1/3
\item -2/3
\item We do not have enough information to know the \((3,1)\) entry of \(\mathbf{L}\).
\end{enumerate}
\smallskip
\textbf{4.} \textit{(choose one)} It is given that

\[\mathbf{A}=\begin{pmatrix}      1   &  0   &  0\\     -1  &   1  &   0\\      0  &   0  &   1 \end{pmatrix}\begin{pmatrix}      1  &  -1 &    1   &  3\\      0 &   -2 &    0   &  4\\      0  &   0 &    3   &  1 \end{pmatrix}.\]

Which of the following is a solution to \(\mathbf{Ax}=\begin{pmatrix}1\\1\\1\end{pmatrix}\)

\begin{enumerate}[label=(\Alph*)]
\item \(\begin{pmatrix}-1\\1\\0\\1\end{pmatrix}\)
\item \(\begin{pmatrix}1\\1\\1\\1\end{pmatrix}\)
\item \(\begin{pmatrix}1\\-1\\1\\1\end{pmatrix}\)
\item \(\begin{pmatrix}1\\-1\\0\\1\end{pmatrix}\)
\end{enumerate}
\smallskip
\textbf{5.} \textit{(choose one)} It is given that \(\mathbf{A}=\begin{pmatrix}        5       &       4      &        5      &        2       \\        1       &       3      &       -1     &         1       \\       -2      &        0     &         3     &        -1    \end{pmatrix}\) has a LU factorization \(\mathbf{A=LU}\), where

\[\mathbf{U}=\begin{pmatrix}        5     &         4          &    5        &      2       \\        0     &        11/5     &      -2    &          3/5     \\        0     &         0          &   71/11  &        -7/11    \end{pmatrix}.\]

What is the \((3,1)\)-entryof \(\mathbf{L}\)?

\begin{enumerate}[label=(\Alph*)]
\item -2/5
\item 5/2
\item -5/2
\item We do not have enough information to know the \((3,1)\) entry of \(\mathbf{L}\).
\item 2/5
\item 1/5
\item -1/5
\end{enumerate}
\smallskip
\section{Quiz 2.8}
\textbf{1.} \textit{(choose one)} Which of the following is \textbf{not} the determinant of \(\begin{pmatrix}1&0&5\\2&3&-1\\1&1&-3\end{pmatrix}\)?

\begin{enumerate}[label=(\Alph*)]
\item \(\begin{vmatrix}3&-1\\1&-3\end{vmatrix}+(2)\begin{vmatrix}0&5\\1&-3\end{vmatrix}+\begin{vmatrix}0&5\\3&-1\end{vmatrix}\)
\item \(\begin{vmatrix}3&-1\\1&-3\end{vmatrix}+5\begin{vmatrix}2&3\\1&1\end{vmatrix}\)
\item \(\begin{vmatrix}3&-1\\1&-3\end{vmatrix}-(2)\begin{vmatrix}0&5\\1&-3\end{vmatrix}+\begin{vmatrix}0&5\\3&-1\end{vmatrix}\)
\item \((5)\begin{vmatrix}2&3\\1&1\end{vmatrix}+\begin{vmatrix}1&0\\1&1\end{vmatrix}-(3)\begin{vmatrix}1&0\\2&3\end{vmatrix}\)
\end{enumerate}
\smallskip
\textbf{2.} \textit{(short answer)} Compute the determinant \(\begin{vmatrix}2&2&1\\0&1&2\\0&-1&1\end{vmatrix}\)

\smallskip
\textbf{3.} \textit{(choose one)} 1.    What \(\begin{vmatrix}1&1&-1\\0&1&1\\0&-1&1\end{vmatrix}\)?

\begin{enumerate}[label=(\Alph*)]
\item 2
\item -2
\item -1
\item 0
\item 1
\end{enumerate}
\smallskip
\textbf{4.} \textit{(choose one)} Which of the following statements is true of the \((3,2)\)-cofactor of \(\begin{pmatrix}a&1&2&3\\2&b&3&4\\3&4&c&5\\4&5&6&1\end{pmatrix}\)?

\begin{enumerate}[label=(\Alph*)]
\item It depends on \(a\).
\item It depends on \(b\).
\item It depends on \(c\).
\item It does not depend on \(a,b,c\).
\end{enumerate}
\smallskip
\textbf{5.} \textit{(choose one)} Which of the following is equivalent to the determinant of \(\begin{pmatrix}1&6&0\\2&-1&3\\1&2&-3\end{pmatrix}\)?

\begin{enumerate}[label=(\Alph*)]
\item \(\left|\begin{array}{cc} -1 & 3 \\ 2 & -3  \end{array}\right|  + 2 \left|\begin{array}{cc} -6 & 0 \\ 2 & -3  \end{array}\right| + \left|\begin{array}{cc} 6 & 0 \\ -1 & 3  \end{array}\right|\)
\item \(\left|\begin{array}{cc}-1 3 \\2 -3\end{array}\right|+ 2 \left|\begin{array}{cc}6 0 \\2 -3\end{array}\right|+ \left|\begin{array}{cc}6 0 \\-1 3\end{array}\right|\)
\item \(\left|\begin{array}{cc} -1 3 \\ 2 -3 \end{array}\right| + 2 \left|\begin{array}{cc} 6 0 \\ -2 -3 \end{array}\right| + \left|\begin{array}{cc} 6 0 \\ -1 3 \end{array}\right|\)
\item \(\left|\begin{array}{cc} -1 & 3 \\ 2 & -3  \end{array}\right| + 6 \left|\begin{array}{cc} -2 & -3 \\ -1 & 3  \end{array}\right|\)
\end{enumerate}
\smallskip
\section{Quiz 2.9}
\textbf{1.} \textit{(short answer)} Suppose \(\mathbf{A}\xrightarrow{R_1+3R_3}\xrightarrow{R_2\leftrightarrow R_1}\xrightarrow{\frac{1}{3}R_3}\mathbf{R}\) and \(\det(\mathbf{R})=2\). What is \(\det(\mathbf{A})\)?

\smallskip
\textbf{2.} \textit{(short answer)} Compute the determinant \(\begin{vmatrix}1&1&2&2&3\\0&1&0&2&0\\0&1&0&2&0\\0&0&0&1&1\\0&0&1&0&0\end{vmatrix}=\).

Note that it can be computed without the use of any calculators.

\smallskip
\textbf{3.} \textit{(short answer)} Compute the determinant \(\begin{vmatrix}1&1&0&0&0\\1&1&1&2&3\\0&0&3&1&-2\\1&1&5&-2&0\\0&0&9&-1&1\end{vmatrix}=\).

Note that it can be computed without the use of any calculators.

\smallskip
\textbf{4.} \textit{(choose one)} Suppose \(\mathbf{A}\xrightarrow{R_2\leftrightarrow R_1}\xrightarrow{3R_3}\xrightarrow{R_1-2R_2}\begin{pmatrix}1&1&3\\0&2&5\\0&0&-1\end{pmatrix}\). What is the determinant of \(\mathbf{A}\)?

\begin{enumerate}[label=(\Alph*)]
\item \(\frac{2}{3}\)
\item \(6\)
\item \(-\frac{2}{3}\)
\item \(-6\)
\item \(0\)
\end{enumerate}
\smallskip
\textbf{5.} \textit{(choose one)} Suppose \(\mathbf{Ax = b}\) has a general solution with 1 parameter. Which of the following is true about the square matrix \(\mathbf{A}\)?

\begin{enumerate}[label=(\Alph*)]
\item The determinant is zero
\item We cannot make any conclusion about the determinant without computing it.
\item The RREF of \(\mathbf{A}\) is a strictly upper triangular matrix.
\item The RREF of \(\mathbf{A}\) is not a triangular matrix.
\end{enumerate}
\smallskip
\section{Quiz 2.10}
\textbf{1.} \textit{(short answer)} Let \(\mathbf{A=LU}\) be a LU-factorization of \(\mathbf{A}\). Suppose \(\det(\mathbf{U})=3\), what is the determinant of \(\mathbf{A}\)?

\smallskip
\textbf{2.} \textit{(choose one)} Which of the following statements regarding the adjoint is false?

\begin{enumerate}[label=(\Alph*)]
\item \(\mathbf{A}adj(\mathbf{A})=det(A)\mathbf{I}\) is true only when \(\mathbf{A}\) is invertible.
\item \(\mathbf{A}adj(\mathbf{A})=det(A)\mathbf{I}\)
\item If \(\mathbf{A}\) is singular, then \( \mathbf{A}adj( \mathbf{A})= \mathbf{0}\).
\end{enumerate}
\smallskip
\textbf{3.} \textit{(choose one)} Let \(\mathbf{A}\) be a \(n\times m\) matrix and \(\mathbf{B}\) a \(m\times n\) matrix. Then \(\det(\mathbf{AB})=\det(\mathbf{A})\det(\mathbf{B})\).

\begin{enumerate}[label=(\Alph*)]
\item False
\item True
\end{enumerate}
\smallskip
\textbf{4.} \textit{(choose one)} Suppose \(\mathbf{A}\) and \(\mathbf{B}\) are square matrix of order 3 with determinants \(\det(\mathbf{A})=5\) and \(\det(\mathbf{B})=2\). What is the determinant of \(3(\mathbf{AB}^T)^{-1}\)?

\begin{enumerate}[label=(\Alph*)]
\item 27/10
\item 3/10
\item 30
\item 270
\item 54/3
\end{enumerate}
\smallskip
\textbf{5.} \textit{(short answer)} Let \(\mathbf{A}=\begin{pmatrix}1&1&2\\0&-1&3\\1&2&-1\end{pmatrix}\). What is the \((3,1)\)-entry of the adjoint matrix \(adj(\mathbf{A})\) of \(\mathbf{A}\)?

\smallskip
\section{Quiz 3.1}
\textbf{1.} \textit{(choose all that apply)} Let \(\mathbf{u}\) be a nonzero vector in \(\mathbb R^n\). Then \(-\frac{1}{2}\mathbf{u}\) is

\begin{enumerate}[label=(\Alph*)]
\item pointing in the opposite direction as \(\mathbf{u}\).
\item shorter than \(\mathbf{u}\).
\item pointing in the same direction as \(\mathbf{u}\).
\item longer than \(\mathbf{u}\).
\end{enumerate}
\smallskip
\textbf{2.} \textit{(choose one)} Let \(a,b\in\mathbb{R}\) be real numbers and \(\mathbf{u,v}\in\mathbb{R}^n\), for \(n>1\). Which statement is true?

\begin{enumerate}[label=(\Alph*)]
\item \(\mathbf{u}^T\mathbf{v}\) is a real number.
\item \(a\mathbf{u}+b\mathbf{v}\) is a real number.
\item \(\mathbf{u+v}\) is a real number.
\item None of the other statements are true.
\end{enumerate}
\smallskip
\textbf{3.} \textit{(short answer)} \(\begin{pmatrix}1\\-1\\0\\0\end{pmatrix}\) is a vector in \(\mathbb{R}^n\), where \(n=\)

\smallskip
\textbf{4.} \textit{(choose one)} Let two vectors \(\mathbf{u}=\begin{pmatrix}1\\0\\2\end{pmatrix}\) and \(\mathbf{v}=\begin{pmatrix}3\\0\\-1\end{pmatrix}\). Which of the following vectors is \textbf{not }a linear combination of \(\mathbf{u}\) and \(\mathbf{v}\)?

\begin{enumerate}[label=(\Alph*)]
\item \(\begin{pmatrix}6\\1\\-2\end{pmatrix}\)
\item \(\begin{pmatrix}4\\0\\-1\end{pmatrix}\)
\item \(\begin{pmatrix}2\\0\\-3\end{pmatrix}\)
\item \(\begin{pmatrix}3\\0\\6\end{pmatrix}\)
\end{enumerate}
\smallskip
\textbf{5.} \textit{(choose one)} \(\begin{pmatrix}1\\1\\0\end{pmatrix}\) is a vector in

\begin{enumerate}[label=(\Alph*)]
\item \(\mathbb{R}^3\)
\item \(\mathbb{R}^2\)
\item \(\mathbb{R}^4\)
\end{enumerate}
\smallskip
\section{Quiz 3.2}
\textbf{1.} \textit{(choose one)} Let \(u,v,w\) be (column) vector in \(\mathbb{R}^n\), for some \(n>1\), and \(\alpha,\beta\) be real numbers. Which statement is FALSE?

\begin{enumerate}[label=(\Alph*)]
\item \(\mathbf{u\cdot v}=\mathbf{u v}^T\)
\item \(\mathbf{u\cdot v=(u\cdot v)}^T\)
\item \(\mathbf{u}\cdot(\alpha \mathbf{v}+\beta \mathbf{w})=\alpha \mathbf{u\cdot v}+\beta \mathbf{u\cdot w}\)
\item \(-1\leq\frac{\mathbf{u\cdot v}}{\Vert\mathbf{u}\Vert\,\Vert \mathbf{v}\Vert}\leq 1\)
\end{enumerate}
\smallskip
\textbf{2.} \textit{(choose one)} What is the norm of \(\begin{pmatrix}1\\1\\-1\end{pmatrix}\)?

\begin{enumerate}[label=(\Alph*)]
\item \(\sqrt{3}\)
\item \(1\)
\item \(\sqrt{2}\)
\item \(0\)
\end{enumerate}
\smallskip
\textbf{3.} \textit{(short answer)} What is \(\begin{pmatrix}1\\1\\1\end{pmatrix}\cdot\begin{pmatrix}1\\0\\2\end{pmatrix}\)?

\smallskip
\textbf{4.} \textit{(choose one)} What is \(\begin{pmatrix}1\\2\\-1\\1\end{pmatrix}\cdot\begin{pmatrix}0\\1\\1\\1\end{pmatrix}\)?

\begin{enumerate}[label=(\Alph*)]
\item 2
\item 1
\item 0
\item 3
\item 4
\end{enumerate}
\smallskip
\textbf{5.} \textit{(choose one)} What can you say about two vectors \(\mathbf{u}\) and \(\mathbf{v}\) in \(\mathbb{R}^n\) if \(\mathbf{u\cdot v}\) is negative?

\begin{enumerate}[label=(\Alph*)]
\item The angle between them is obtuse.
\item One of the vectors is a multiple of the other.
\item The angle between them is acute.
\item They are perpendicular (the angle between them is \(\pi/2\))
\end{enumerate}
\smallskip
\section{Quiz 3.3}
\textbf{1.} \textit{(choose one)} Let \(S=\{\mathbf{u}_1,...,\mathbf{u}_k\}\subseteq\mathbb{R}^n\) and \(\mathbf{A}=\begin{pmatrix}\mathbf{u}_1&\cdots&\mathbf{u}_k\end{pmatrix}\). Which set \(S\) does not spans \(\mathbb{R}^n\)?

\begin{enumerate}[label=(\Alph*)]
\item \(\mathbf{A}\rightarrow\begin{pmatrix}1&0&0&2\\0&1&1&1\\0&0&0&0\end{pmatrix}\)
\item \(\mathbf{A}\rightarrow\begin{pmatrix}1&0&1&0\\0&1&-1&0\\0&0&1&0\end{pmatrix}\)
\item \(\mathbf{A}\rightarrow\begin{pmatrix}1&0&0&1\\0&1&0&-1\\0&0&1&0\end{pmatrix}\)
\item \(\mathbf{A}\rightarrow\begin{pmatrix}1&0&0\\0&1&0\\0&0&1\end{pmatrix}\)
\end{enumerate}
\smallskip
\textbf{2.} \textit{(choose all that apply)} Let \(\{\mathbf{u}_1,\mathbf{u}_2,\mathbf{u}_3,\mathbf{u}_4\}\) be a set of vectors in \(\mathbb{R}^3\). Suppose the reduced row-echelon form of \(\mathbf{A}=\begin{pmatrix}\mathbf{u}_1&\mathbf{u}_2&\mathbf{u}_3&\mathbf{u}_4\end{pmatrix}\) is \(\begin{pmatrix}1&0&0&0\\0&1&0&-1\\0&0&1&1\end{pmatrix}\). Which of the following statements is true?

\begin{enumerate}[label=(\Alph*)]
\item \(\{\mathbf{u}_1,\mathbf{u}_2,\mathbf{u}_4\}\) spans \(\mathbb{R}^3\)
\item \(\{\mathbf{u}_1,\mathbf{u}_2,\mathbf{u}_3\}\) spans \(\mathbb{R}^3\)
\item \(\mathbf{u}_2\) is a linear combination of \(\mathbf{u}_3\) and \(\mathbf{u}_4\).
\item \(\{\mathbf{u}_2,\mathbf{u}_3,\mathbf{u}_4\}\) spans \(\mathbb{R}^3\)
\end{enumerate}
\smallskip
\textbf{3.} \textit{(choose one)} \(span\left\{\begin{pmatrix}1\\0\\-1\end{pmatrix},\begin{pmatrix}1\\1\\1\end{pmatrix},\begin{pmatrix}1\\0\\0\end{pmatrix},\begin{pmatrix}0\\0\\1\end{pmatrix},\begin{pmatrix}0\\1\\1\end{pmatrix}\right\}\) is a subset of

\begin{enumerate}[label=(\Alph*)]
\item \(\mathbb{R}^3\)
\item \(\mathbb{R}^2\)
\item \(\mathbb{R}^4\)
\item \(\mathbb{R}^5\)
\end{enumerate}
\smallskip
\textbf{4.} \textit{(choose one)} Does \(\begin{pmatrix}1\\0\\0\end{pmatrix}\), \(\begin{pmatrix}1\\1\\0\end{pmatrix}\), \(\begin{pmatrix}1\\1\\1\end{pmatrix}\) span \(\mathbb{R}^3\)?

\begin{enumerate}[label=(\Alph*)]
\item Yes
\item No
\end{enumerate}
\smallskip
\textbf{5.} \textit{(choose one)} Let \(S=\{\mathbf{u_1,u_2,...,u_k}\}\) and \(T=\{\mathbf{v_1,v_2,...,v_m}\}\). To check if \(span(S)\subseteq span(T)\), we check for consistency of

\begin{enumerate}[label=(\Alph*)]
\item \(\left(\begin{array}{cccc|cccc}\mathbf{v_1}&\mathbf{v_2}&\cdots&\mathbf{v_m}&\mathbf{u_1}&\mathbf{u_2}&\cdots&\mathbf{u_k}\end{array}\right)\)
\item \(\left(\begin{array}{cccc|cccc}\mathbf{u_1}&\mathbf{u_2}&\cdots&\mathbf{u_k}&\mathbf{v_1}&\mathbf{v_2}&\cdots&\mathbf{v_m}\end{array}\right)\)
\end{enumerate}
\smallskip
\section{Quiz 3.4}
\textbf{1.} \textit{(choose all that apply)} Which of the following sets are subspace?

\begin{enumerate}[label=(\Alph*)]
\item \(\left\{\begin{array}{c|c}\begin{pmatrix}x\\y\\z\end{pmatrix}&z=x-2y\end{array}\right\}\)
\item \(\left\{\begin{array}{c|c}\begin{pmatrix}x\\y\end{pmatrix}&y=x^2\end{array}\right\}\)
\item \(\left\{\begin{array}{c|c}\begin{pmatrix}x\\y\end{pmatrix}&y=2x-3\end{array}\right\}\)
\item \(\left\{\begin{array}{c|c}\begin{pmatrix}a\\b\\c\end{pmatrix}&abc=0\end{array}\right\}\)
\end{enumerate}
\smallskip
\textbf{2.} \textit{(choose one)} Which of the following statement is FALSE?

\begin{enumerate}[label=(\Alph*)]
\item Every subspace contains infinitely many vectors.
\item Every subspace must contain the zero vector.
\item The set containing only the zero vector is a subspace.
\item Every nonzero subspace must contain a line.
\end{enumerate}
\smallskip
\textbf{3.} \textit{(choose one)} Which set is a subspace?

\begin{enumerate}[label=(\Alph*)]
\item \(\left\{\begin{array}{c|c}s\begin{pmatrix}1\\0\\1\end{pmatrix}+t\begin{pmatrix}0\\-1\\0\end{pmatrix}&s,t\in\mathbb{R}\end{array}\right\}\)
\item \(\left\{\begin{array}{c|c}\begin{pmatrix}s\\t\\1\end{pmatrix}&s,t\in\mathbb{R}\end{array}\right\}\)
\item \(\left\{\begin{array}{c|c}\begin{pmatrix}s\\s^2\\0\end{pmatrix}&s\in\mathbb{R}\end{array}\right\}\)
\end{enumerate}
\smallskip
\textbf{4.} \textit{(choose one)} Which of the following is not a criteria in the definition for a subset \(V\subseteq\mathbb R^n\) to be a subspace?

\begin{enumerate}[label=(\Alph*)]
\item There exists a finite set \(S=\{\mathbf{u}_1,\mathbf{u}_2,...,\mathbf{u}_k\}\subseteq\mathbb{R}^n\) such that \(V=\text{span}(S)\).
\item Contains the origin \(\mathbf{0}\in V\).
\item For any vector \(\mathbf{v}\in V\) and real number \(\alpha\in\mathbb{R}\), \(\alpha\mathbf{v}\in V\).
\item For any vectors \(\mathbf{u},\mathbf{v}\in V\), \(\mathbf{u}+\mathbf{v}\in V\).
\end{enumerate}
\smallskip
\textbf{5.} \textit{(choose one)} Is the set \(V=\{(x,y,z)\mid 2x+3y-z=1,\;x-3z=2\}\) a subspace?

\begin{enumerate}[label=(\Alph*)]
\item No
\item Yes
\end{enumerate}
\smallskip
\section{Quiz 3.5}
\textbf{1.} \textit{(choose all that apply)} Suppose \(\{\mathbf{u}_1,\mathbf{u}_2,\mathbf{u}_3\}\) is a linearly independent set. Let \(\mathbf{v}\)  be a vector such that \((\begin{array}{ccc|c}\mathbf{u}_1&\mathbf{u}_2&\mathbf{u}_3&\mathbf{v}\end{array})\longrightarrow\left(\begin{array}{ccc|c}1&0&0&-2\\0&1&0&0\\0&0&1&1\\0&0&0&1\end{array}\right)\). Which of the following statements are true?

\begin{enumerate}[label=(\Alph*)]
\item \(\{​​\mathbf{u}_1,\mathbf{u}_2,\mathbf{u}_3,\mathbf{v}\}\) is a linearly independent set
\item \(\mathbf{v}\)  belongs to \(\text{span}\{​​\mathbf{u}_1,\mathbf{u}_2,\mathbf{u}_3\}\).
\item \(​\mathbf{v}\) is a linear combination of \(\mathbf{u}_1\) and \(\mathbf{u}_3\).
\item \(\{​​\mathbf{u}_1,\mathbf{u}_2,\mathbf{v}\}\) is a linearly independent set.
\end{enumerate}
\smallskip
\textbf{2.} \textit{(choose all that apply)} Suppose the set \(\{\mathbf{u},\mathbf{v},\mathbf{w}\}\) in \(\mathbb{R}^3\) is linearly dependent with \(\mathbf{u}\neq\mathbf{0}\). What are the possible geometrical object of \(\text{span}(\{\mathbf{u},\mathbf{v},\mathbf{w}\})\)?

\begin{enumerate}[label=(\Alph*)]
\item a line
\item a point
\item a plane
\item the whole \(\mathbb{R}^3\)
\end{enumerate}
\smallskip
\textbf{3.} \textit{(choose all that apply)} Which of the following sets are linearly independent?

\begin{enumerate}[label=(\Alph*)]
\item \(\left\{\begin{pmatrix}1\\1\end{pmatrix}\right\}\)
\item \(\left\{\begin{pmatrix}1\\1\end{pmatrix},\begin{pmatrix}0\\0\end{pmatrix}\right\}\)
\item \(\left\{\begin{pmatrix}1\\2\end{pmatrix},\begin{pmatrix}1\\1\end{pmatrix}\right\}\)
\item \(\left\{\begin{pmatrix}1\\1\\0\end{pmatrix},\begin{pmatrix}1\\2\\1\end{pmatrix}\right\}\)
\item \(\left\{\begin{pmatrix}1\\1\\0\end{pmatrix},\begin{pmatrix}0\\1\\1\end{pmatrix},\begin{pmatrix}1\\2\\1\end{pmatrix}\right\}\)
\end{enumerate}
\smallskip
\textbf{4.} \textit{(choose one)} Which of the following statements about linear independence is true?

\begin{enumerate}[label=(\Alph*)]
\item A set with more than \(n\) vectors in \(\mathbb{R}^n\) must be linearly dependent.
\item A set with less than \(n\) vectors in \(\mathbb{R}^n\) must be linearly independent
\item A set with exactly \(n\) vectors in \(\mathbb{R}^n\) must be linearly independent .
\item A set with exactly \(n\) vectors in \(\mathbb{R}^n\) such that none of the vectors in the set is a multiple of another must be linearly independent.
\end{enumerate}
\smallskip
\textbf{5.} \textit{(true or false)} \(\text{span}\left\{\begin{pmatrix}1\\1\\0\end{pmatrix},\begin{pmatrix}1\\2\\1\end{pmatrix},\begin{pmatrix}1\\0\\-1\end{pmatrix}\right\}=\text{span}\left\{\begin{pmatrix}1\\1\\0\end{pmatrix},\begin{pmatrix}1\\2\\1\end{pmatrix}\right\}\)

\smallskip
\section{Quiz 3.6}
\textbf{1.} \textit{(choose one)} Suppose the rows of matrix \(\mathbf{A}\) is a basis for \(\mathbb{R}^3\). Which statement is FALSE?

\begin{enumerate}[label=(\Alph*)]
\item \(\mathbf{A}\) is not necessary a square matrix.
\item The columns of \(\mathbf{A}\) must be linearly independent.
\item \(\mathbf{Ax=0}\) has only the trivial solution.
\item \(\mathbf{A}\) is a 3 by 3 invertible matrix.
\end{enumerate}
\smallskip
\textbf{2.} \textit{(choose one)} Let \(S=\{u_1,u_2,u_3\}\) be a subset of a subspace \(V\), and \(\mathbf{v} \) a vector in \(V\). Suppose \(\left(\begin{array}{ccc|c}\mathbf{u}_1&\mathbf{u}_2&\mathbf{u}_3&\mathbf{v}\end{array}\right)\rightarrow\left(\begin{array}{ccc|c}1&0&-1&2\\0&1&-1&1\\0&0&0&0\end{array}\right)\). What can we conclude?

\begin{enumerate}[label=(\Alph*)]
\item \(S\) is not linearly independent
\item \(V\subseteq span(S)\)
\item \(span(S)=V\)
\item \(S\) is a basis for \(V\)
\end{enumerate}
\smallskip
\textbf{3.} \textit{(choose all that apply)} Which of the following statement is true about the zero space \(\{\mathbf{0}\}\)?

\begin{enumerate}[label=(\Alph*)]
\item It is a subspace
\item The zero space is the span of \(\{\mathbf{0}\}\).
\item \(\{\mathbf{0}\}\) is a basis for the zero space.
\item \(\{\mathbf{0}\}\) is linearly independent since it contains only 1 vector.
\end{enumerate}
\smallskip
\textbf{4.} \textit{(choose one)} Suppose \(V=\text{span}\{\mathbf{u}_1,\mathbf{u}_2,\mathbf{u}_3,\mathbf{u}_4\}\), where \(\mathbf{u}_1,\mathbf{u}_2,\mathbf{u}_3,\mathbf{u}_4\in\mathbb{R}^4\) and the RREF of \(\left(\begin{array}{cccc}\mathbf{u}_1&\mathbf{u}_2&\mathbf{u}_3&\mathbf{u}_4\end{array}\right)\) is \(\begin{pmatrix}1&0&2&0\\0&1&-1&0\\0&0&0&1\\0&0&0&0\end{pmatrix}\). Which of the following set is not a basis for \(V\)?

\begin{enumerate}[label=(\Alph*)]
\item \(\{\mathbf{u}_1,\mathbf{u}_2,\mathbf{u}_3\}\)
\item \(\{\mathbf{u}_2,\mathbf{u}_3,\mathbf{u}_4\}\)
\item \(\{\mathbf{u}_1,\mathbf{u}_2,\mathbf{u}_4\}\)
\item \(\{\mathbf{u}_1,\mathbf{u}_3,\mathbf{u}_4\}\)
\end{enumerate}
\smallskip
\textbf{5.} \textit{(choose one)} Let \(S=\{\mathbf{u}_1,\mathbf{u}_2,\mathbf{u}_3\}\) be a basis for \(V\), a subspace of \(\mathbb{R}^n\). Which of the following statements is \textbf{false}?

\begin{enumerate}[label=(\Alph*)]
\item If \(n3\), then there is a vector \(\mathbf{v}\) in \(V\) such that \([\mathbf{v}]_S\) is not well-defined.
\item If \(n=3\), then for every vector \(\mathbf{v}\) in \(\mathbb{R}^n\), \([\mathbf{v}]_S\) is well-defined.
\item If \(n3\), then there is a vector \(\mathbf{v}\) in \(\mathbb{R}^n\) such that \([\mathbf{v}]_S\) is not well-defined.
\item If \(n=3\), then for every vector \(\mathbf{v}\) in \(V\), \([\mathbf{v}]_S\) is well-defined.
\end{enumerate}
\smallskip
\section{Quiz 3.7}
\textbf{1.} \textit{(choose one)} Let \(V\) be a subspace of dimension 5. Which of the following statement is \textbf{false}?

\begin{enumerate}[label=(\Alph*)]
\item Any set with 5 vectors will span \(V\).
\item There are 5 vectors in any basis of \(V\).
\item Any set containing 7 vectors in \(V\) must be linearly dependent.
\item Any set that spans \(V\) must contain at least 5 vectors.
\end{enumerate}
\smallskip
\textbf{2.} \textit{(choose one)} Suppose \(V=\text{span}\{\mathbf{u}_1.\mathbf{u}_2,\mathbf{u}_3\}\). Let \(\mathbf{v}\in V\) be a vector such that \(\mathbf{v}=3\mathbf{u}_1-2\mathbf{u}_2-\mathbf{u}_3\) and \(\mathbf{v}=3\mathbf{u}_2+\mathbf{u}_3\). Which of the following statement must be true?

\begin{enumerate}[label=(\Alph*)]
\item Any basis for \(V\) must contain less than 3 vectors.
\item \(\{\mathbf{u}_1.\mathbf{u}_2,\mathbf{u}_3\}\) is a basis for \(V\).
\item Any basis for \(V\) must contain 3 vectors.
\item Any subset of \(\{\mathbf{u}_1.\mathbf{u}_2,\mathbf{u}_3\}\) containing 2 vectors is a basis for \(V\).
\end{enumerate}
\smallskip
\textbf{3.} \textit{(choose all that apply)} Let \(U\) and \(V\) be subspaces of \(\mathbb{R}^5\) such that \(U\subseteq V\). Suppose \(\{\mathbf{u}_1,\mathbf{u}_2,\mathbf{u}_3\}\) is a basis for a subspace \(U\) and \(\{\mathbf{v}_1,\mathbf{v}_2,\mathbf{v}_3,\mathbf{v}_4\}\) is a linearly dependent set that spans \(V\). Which of the following statements are true?

\begin{enumerate}[label=(\Alph*)]
\item Dimension of \(U\) is 3.
\item \(U=V\)
\item Dimension of \(V\) is 4.
\item \(\{\mathbf{v}_1,\mathbf{v}_2,\mathbf{v}_3,\mathbf{v}_4\}\) is a basis for \(V\).
\end{enumerate}
\smallskip
\textbf{4.} \textit{(choose one)} Let \(U\) be a subspace of \(\mathbb{R}^n\). Suppose \(V\) is a 3-dimensional subspace of  \(\mathbb{R}^n\). Which of the following statements is false?

\begin{enumerate}[label=(\Alph*)]
\item If \(U\) is a 3-dimensional subspace of \(\mathbb R^n\), then necessarily \(U=V\).
\item If \(U\subseteq V\) and \(\{\mathbf{u}_1,\mathbf{u}_2,\mathbf{u}_3\}\) is a linearly independent subset of \(U\), then \(\{\mathbf{u}_1,\mathbf{u}_2,\mathbf{u}_3\}\) is a basis for \(U\).
\item If \(span\{\mathbf{u}_1,\mathbf{u}_2\}=U\), then \(U\neq V\).
\item If \(V\subseteq U\), then \(\{\mathbf{u}_1,\mathbf{u}_2\}\) cannot span \(U\).
\end{enumerate}
\smallskip
\textbf{5.} \textit{(choose one)} Let \(V=\left\{\begin{array}{c|c}\begin{pmatrix}x\\y\\z\end{pmatrix}&3x-2y+z=0\end{array}\right\}\). Which of the following sets form a basis for \(V\)?

\begin{enumerate}[label=(\Alph*)]
\item \(\left\{\begin{pmatrix}2\\3\\0\end{pmatrix},\begin{pmatrix}1\\0\\-3\end{pmatrix}\right\}\)
\item \(\left\{\begin{pmatrix}3\\3\\-3\end{pmatrix},\begin{pmatrix}1\\1\\-1\end{pmatrix}\right\}\)
\item \(\left\{\begin{pmatrix}2\\3\\0\end{pmatrix},\begin{pmatrix}1\\0\\-3\end{pmatrix},\begin{pmatrix}1\\1\\1\end{pmatrix}\right\}\)
\item \(\left\{\begin{pmatrix}1\\1\\1\end{pmatrix}\right\}\)
\end{enumerate}
\smallskip
\section{Quiz 3.8}
\textbf{1.} \textit{(choose one)} Let \(S=\{\mathbf{u}_1,\mathbf{u}_2,\mathbf{u}_3\}\) and \(T=\{\mathbf{v}_1,\mathbf{v}_2,\mathbf{v}_3\}\) be 2 bases for a subspace \(V\subseteq\mathbb{R}^9\).

Suppose \(\mathbf{P}=\begin{pmatrix}      0  &   1  &   1\\      0   &  2 &    1\\      1   &  0 &   -1 \end{pmatrix}\) is the transition matrix from \(S\) to \(T\). Let \(\mathbf{w}\) be a vector in \(V\) such that \([\mathbf{w}]_S=\begin{pmatrix}1\\2\\-1\end{pmatrix}\). 

What is \([\mathbf{w}]_T\)?

\begin{enumerate}[label=(\Alph*)]
\item \([\mathbf{w}]_T=\begin{pmatrix}1\\3\\2\end{pmatrix}\)
\item \([\mathbf{w}]_T=\begin{pmatrix}-1\\1\\0\end{pmatrix}\)
\item \([\mathbf{w}]_T=\begin{pmatrix}-1\\-2\\1\end{pmatrix}\)
\item \([\mathbf{w}]_T=\begin{pmatrix}1\\2\\3\end{pmatrix}\)
\end{enumerate}
\smallskip
\textbf{2.} \textit{(choose one)} Let \(S=\{\mathbf{u}_1,\mathbf{u}_2,\mathbf{u}_3\}\) and \(T=\{\mathbf{v}_1,\mathbf{v}_2,\mathbf{v}_3\}\) be 2 bases for \(\mathbb{R}^3\). Suppose \((\begin{array}{ccc|ccc}\mathbf{u}_1&\mathbf{u}_2&\mathbf{u}_3&\mathbf{v}_1&\mathbf{v}_2&\mathbf{v}_3\end{array})\xrightarrow{RREF}(\begin{array}{c|c}\mathbf{I}&\mathbf{P}\end{array})\). Which of the following statement is true?

\begin{enumerate}[label=(\Alph*)]
\item \(\mathbf{P}\) is the transition matrix from basis \(T\) to \(S\).
\item \(\mathbf{P}\) is the transition matrix from basis \(S\) to \(T\).
\item None of the other options are true.
\end{enumerate}
\smallskip
\textbf{3.} \textit{(choose one)} Let \(S=\{\mathbf{v}_1,\mathbf{v}_2,\mathbf{v}_3\}\) and \(T=\{\mathbf{u}_1,\mathbf{u}_2,\mathbf{u}_3\}\) be two bases of a subspace \(V\subseteq\mathbb{R}^5\). Suppose the RREF of \((\begin{array}{ccc|c|c|c}\mathbf{u}_1&\mathbf{u}_2&\mathbf{u}_3&\mathbf{v}_1&\mathbf{v}_2&\mathbf{v}_3\end{array})\) is \(\left(\begin{array}{ccc|c|c|c}1&0&0&2&-1&-1\\0&1&0&0&3&1\\0&0&1&1&1&3\\0&0&0&0&0&0\\0&0&0&0&0&0\end{array}\right)\). Which of the following statements is true?

\begin{enumerate}[label=(\Alph*)]
\item Let \(\mathbf{P}=\begin{pmatrix}2&-1&-1\\0&3&1\\1&1&3\end{pmatrix}\). For any vector \(\mathbf{w}\in V\), \([\mathbf{w}]_S=\mathbf{P}^{-1}[\mathbf{w}]_T\).
\item Let \(\mathbf{P}=\begin{pmatrix}2&-1&-1\\0&3&1\\1&1&3\end{pmatrix}\). For any vector \(\mathbf{w}\in V\), \([\mathbf{w}]_S=\mathbf{P}[\mathbf{w}]_T\).
\item Let \(\mathbf{P}=\begin{pmatrix}2&-1&-1\\0&3&1\\1&1&3\\0&0&0\\0&0&0\end{pmatrix}\). For any vector \(\mathbf{w}\in V\), \([\mathbf{w}]_S=\mathbf{P}[\mathbf{w}]_T\).
\item Let \(\mathbf{P}=\begin{pmatrix}2&-1&-1\\0&3&1\\1&1&3\\0&0&0\\0&0&0\end{pmatrix}\). For any vector \(\mathbf{w}\in V\), \([\mathbf{w}]_T=\mathbf{P}[\mathbf{w}]_S\).
\item Let \(\mathbf{P}=\begin{pmatrix}2&-1&-1\\0&3&1\\1&1&3\\0&0&0\\0&0&0\end{pmatrix}\). For any vector \(\mathbf{w}\in V\), \([\mathbf{w}]_S=\mathbf{P}[\mathbf{w}]\).
\item Let \(\mathbf{P}=\begin{pmatrix}2&-1&-1\\0&3&1\\1&1&3\\0&0&0\\0&0&0\end{pmatrix}\). For any vector \(\mathbf{w}\in V\), \([\mathbf{w}]_T=\mathbf{P}[\mathbf{w}]\).
\end{enumerate}
\smallskip
\textbf{4.} \textit{(choose one)} Let \(S\) and \(T\) be 2 bases for a subspace \(V\subseteq\mathbb{R}^n\). Suppose \(\mathbf{P}\) is the transition matrix from \(S\) to \(T\) and \( (\begin{array}{c|c} \mathbf{P} & \mathbf{I} \end{array}) \rightarrow ( \begin{array}{c|c}\mathbf{I}&\mathbf{Q} \end{array} ) \). Then \(\mathbf{Q}\) is the transition matrix from \(T\) to \(S\).

\begin{enumerate}[label=(\Alph*)]
\item True
\item False
\end{enumerate}
\smallskip
\textbf{5.} \textit{(choose one)} Let \(S\) and \(T\) be 2 bases for a subspace \(V\subseteq\mathbb{R}^n\).

Suppose \(\mathbf{P}=\begin{pmatrix}      0  &   1  &   1\\      0   &  2 &    1\\      1   &  0 &   -1 \end{pmatrix}\) is the transition matrix from \(S\) to \(T\). Which of the following is the transition matrix from \(T\) to \(S\)?

\begin{enumerate}[label=(\Alph*)]
\item \(\begin{pmatrix} 2 -1 1\\ -1 1 0\\ 2 -1 0\end{pmatrix}\)
\item \(\begin{pmatrix} -1 0 2\\ 0 1 0\\ 2 -1 1\end{pmatrix}\)
\item \(\begin{pmatrix} 1 1 0\\ 2 -1 2\\ 2 -1 0\end{pmatrix}\)
\item \(\begin{pmatrix} 2 0 1\\ -1 -1 0\\ 2 0 0\end{pmatrix}\)
\end{enumerate}
\smallskip
\section{Quiz 4.1}
\textbf{1.} \textit{(choose all that apply)} Which vector(s) is/are in the nullspace of \(\begin{pmatrix}1&1&-1&0\\2&1&0&1\\0&-1&2&-1\end{pmatrix}\)?

\begin{enumerate}[label=(\Alph*)]
\item \(\begin{pmatrix}0\\0\\0\\0\end{pmatrix}\)
\item \(\begin{pmatrix}-1\\2\\1\\0\end{pmatrix}\)
\item \(\begin{pmatrix}0\\0\\0\end{pmatrix}\)
\item \(\begin{pmatrix}-1\\2\\1\end{pmatrix}\)
\end{enumerate}
\smallskip
\textbf{2.} \textit{(choose one)} Consider the following matrix

\[\mathbf{A}=\begin{pmatrix}1&2&3\\4&5&6\\7&8&9\end{pmatrix}.\]

Which of the following sets of vectors is a basis for the column space of \(\mathbf{A}\)?

Trivia: The column space is the same as its row space.

\begin{enumerate}[label=(\Alph*)]
\item \(\left\{\begin{pmatrix}1\\0\\-1\end{pmatrix},\begin{pmatrix}0\\1\\2\end{pmatrix}\right\}\)
\item \(\left\{\begin{pmatrix}1\\0\\0\end{pmatrix},\begin{pmatrix}0\\1\\0\end{pmatrix},\begin{pmatrix}0\\0\\1\end{pmatrix}\right\}\)
\item \(\left\{\begin{pmatrix}1\\4\\7\end{pmatrix},\begin{pmatrix}2\\5\\8\end{pmatrix},\begin{pmatrix}1\\2\\3\end{pmatrix}\right\}\)
\item \(\left\{\begin{pmatrix}1\\0\\0\end{pmatrix},\begin{pmatrix}0\\1\\0\end{pmatrix}\right\}\)
\end{enumerate}
\smallskip
\textbf{3.} \textit{(choose all that apply)} Let \(\mathbf{A}\) be a \(5\times 7\) matrix. Which of the following statements are true?

\begin{enumerate}[label=(\Alph*)]
\item The row space of \(\mathbf{A}\) is a subspace of \(\mathbb{R}^5\).
\item The row space of \(\mathbf{A}\) is a subspace of \(\mathbb{R}^7\).
\item The column space of \(\mathbf{A}\) is a subspace of \(\mathbb{R}^5\).
\item The column space of \(\mathbf{A}\) is a subspace of \(\mathbb{R}^7\).
\end{enumerate}
\smallskip
\textbf{4.} \textit{(choose one)} Let \(\mathbf{A}\) be a \(m\times n\) matrix and \(\mathbf{R}\) its reduced row-echelon form. Then the

\begin{enumerate}[label=(\Alph*)]
\item nonzero rows of \(\mathbf{R}\) form a basis for the row space of \(\mathbf{A}\).
\item nonzero rows of \(\mathbf{A}\) form a basis for the row space of \(\mathbf{A}\).
\item rows of \(\mathbf{A}\) has the same linear relations as the rows of \(\mathbf{R}\).
\end{enumerate}
\smallskip
\textbf{5.} \textit{(choose one)} Which of the following is in the nullspace of \(\mathbf{A}=\begin{pmatrix}      1  &   0  &  -1  &   1 &  -1\\      4  &   0  &  -4  &   2 &    0\\      1  &   0  &  -1  &   0  &   1 \end{pmatrix}\)?

\begin{enumerate}[label=(\Alph*)]
\item \(\begin{pmatrix}0\\1\\1\\2\\1\end{pmatrix}\)
\item \(\begin{pmatrix}1\\0\\1\\2\\1\end{pmatrix}\)
\item \(\begin{pmatrix}1\\1\\0\\2\\1\end{pmatrix}\)
\item \(\begin{pmatrix}1\\1\\2\\0\\1\end{pmatrix}\)
\item \(\begin{pmatrix}1\\1\\2\\1\\0\end{pmatrix}\)
\end{enumerate}
\smallskip
\section{Quiz 4.2}
\textbf{1.} \textit{(choose all that apply)} Suppose \(\begin{vmatrix}1&1&0&1&0\\2&1&0&1&0\\1&2&0&1&1\\1&1&0&2&1\\0&0&2&0&2\end{vmatrix}=4\). Which of the following statements are true?

\begin{enumerate}[label=(\Alph*)]
\item \(\left\{\begin{pmatrix}1\\2\\1\\1\\0\end{pmatrix},\begin{pmatrix}1\\1\\2\\1\\0\end{pmatrix},\begin{pmatrix}0\\0\\0\\0\\2\end{pmatrix},\begin{pmatrix}1\\1\\1\\2\\0\end{pmatrix},\begin{pmatrix}0\\0\\1\\1\\2\end{pmatrix}\right\}\) form a basis for \(\mathbb{R}^5\).
\item \(\left\{\begin{pmatrix}2&1&0&1&0\end{pmatrix},\begin{pmatrix}1&1&0&1&0\end{pmatrix},\begin{pmatrix}1&1&0&2&1\end{pmatrix}\right\}\) is linearly independent.
\item The rows form a basis for \(\mathbb{R}^5\).
\end{enumerate}
\smallskip
\textbf{2.} \textit{(true or false)} Suppose \(\mathbf{A}\) is an order 5 square matrix and \(\begin{pmatrix}1\\1\\0\\0\\1\end{pmatrix}\) is a solution to \(\mathbf{Ax}=\mathbf{0}\). Then the rank of \(\mathbf{A}\) must be less than or equal to 4.

\smallskip
\textbf{3.} \textit{(true or false)} Suppose \(\mathbf{A}\) has full rank. Then \(nullity(A)=0\).

\smallskip
\textbf{4.} \textit{(true or false)} Suppose \(\mathbf{A}\) is an order 3 square matrix such that \(\mathbf{v}=\begin{pmatrix}1\\-1\\0\end{pmatrix}\) is a solution to \(\mathbf{Ax=0}\). Then it must be that the columns of \(\mathbf{A}\) cannot be a basis for \(\mathbb{R}^3\).

\smallskip
\textbf{5.} \textit{(choose all that apply)} Given that a matrix \(\mathbf{A}\) is row equivalent to \(\begin{pmatrix}1&0&2&1&0&2\\0&0&0&3&2&0\\0&0&0&0&5&1\end{pmatrix}\). Which of the following statements are true?

\begin{enumerate}[label=(\Alph*)]
\item \(\mathbf{A}\) is of full rank.
\item The rank of \(\mathbf{A}\) is strictly lesser than the number of columns.
\item The dimension of the columns space of \(\mathbf{A}\) is more than 3.
\item The dimension of the row space of \(\mathbf{A}\) is more than 3.
\end{enumerate}
\smallskip
\section{Quiz 5.1}
\textbf{1.} \textit{(choose one)} Which of the following pairs of vectors is not orthogonal?

\begin{enumerate}[label=(\Alph*)]
\item \(\begin{pmatrix}2\\3\end{pmatrix},\begin{pmatrix}-1\\1\end{pmatrix}\)
\item \(\begin{pmatrix}1\\2\\3\end{pmatrix},\begin{pmatrix}0\\3\\-2\end{pmatrix}\)
\item \(\begin{pmatrix}0\\0\\0\end{pmatrix},\begin{pmatrix}1\\2\\1\end{pmatrix}\)
\item \(\begin{pmatrix}1\\-1\\1\\-1\end{pmatrix},\begin{pmatrix}2\\2\\2\\2\end{pmatrix}\)
\end{enumerate}
\smallskip
\textbf{2.} \textit{(choose one)} Let \(V=\{(x,y,z)\mid 2x-y+3z=0\}\). Which vector is orthogonal to \(V\)?

\begin{enumerate}[label=(\Alph*)]
\item \(\begin{pmatrix}-4\\2\\-6\end{pmatrix}\)
\item \(\begin{pmatrix}1\\1\\1\end{pmatrix}\)
\item \(\begin{pmatrix}1\\2\\0\end{pmatrix}\)
\item \(\begin{pmatrix}1\\-1\\-1\end{pmatrix}\)
\end{enumerate}
\smallskip
\textbf{3.} \textit{(choose one)} Let \(V=span\{\begin{pmatrix}1\\0\\-1\end{pmatrix},\begin{pmatrix}1\\0\\1\end{pmatrix}\}\). Which vector is orthogonal to \(V\)?

\begin{enumerate}[label=(\Alph*)]
\item \(\begin{pmatrix}0\\-2\\0\end{pmatrix}\)
\item \(\begin{pmatrix}1\\1\\1\end{pmatrix}\)
\item \(\begin{pmatrix}1\\-2\\1\end{pmatrix}\)
\item \(\begin{pmatrix}-1\\2\\1\end{pmatrix}\)
\end{enumerate}
\smallskip
\textbf{4.} \textit{(choose all that apply)} Suppose \(V=\left\{\begin{array}{c|c}\begin{pmatrix}x\\y\\z\end{pmatrix}&2x-y+3z=0\end{array}\right\}\). Which vector is \textbf{not }orthogonal to \(V\)?

\begin{enumerate}[label=(\Alph*)]
\item \(\begin{pmatrix}1\\1\\1\end{pmatrix}\)
\item \(\begin{pmatrix}0\\0\\0\end{pmatrix}\)
\item \(\begin{pmatrix}2\\-1\\3\end{pmatrix}\)
\item \(\begin{pmatrix}4\\-2\\6\end{pmatrix}\)
\end{enumerate}
\smallskip
\textbf{5.} \textit{(choose all that apply)} Suppose \(\left\{\begin{pmatrix}1\\1\\0\\0\end{pmatrix},\begin{pmatrix}2\\0\\0\\1\end{pmatrix},\begin{pmatrix}1\\0\\1\\0\end{pmatrix}\right\}\) spans a subspace \(V\). Which of the following vectors are orthogonal to \(V\)?

\begin{enumerate}[label=(\Alph*)]
\item \(\begin{pmatrix}0\\0\\0\\0\end{pmatrix}\)
\item \(\begin{pmatrix}-1\\1\\1\\2\end{pmatrix}\)
\item \(\begin{pmatrix}-9\\9\\9\\18\end{pmatrix}\)
\item \(\begin{pmatrix}-1\\1\\1\\1\end{pmatrix}\)
\item \(\begin{pmatrix}-1\\1\\1\\0\end{pmatrix}\)
\end{enumerate}
\smallskip
\section{Quiz 5.2}
\textbf{1.} \textit{(choose one)} Let \(W\) be a subspace of dimension 3. We can never find an orthonormal subset of \(W\) containing 4 vectors.

\begin{enumerate}[label=(\Alph*)]
\item True
\item False
\end{enumerate}
\smallskip
\textbf{2.} \textit{(choose one)} Let \(S=\left\{\frac{1}{\sqrt{2}}\begin{pmatrix}1\\-1\\0\end{pmatrix},\frac{1}{\sqrt{2}}\begin{pmatrix}1\\1\\0\end{pmatrix},\begin{pmatrix}0\\0\\1\end{pmatrix}\right\}\) be an orthonormal basis for \(\mathbb{R}^3\). Let \(\mathbf{w}=\begin{pmatrix}2\\1\\3\end{pmatrix}\). What is \([\mathbf{w}]_S\)?

\begin{enumerate}[label=(\Alph*)]
\item \(\begin{pmatrix}\frac{1}{\sqrt{2}}\\\frac{3}{\sqrt{2}}\\3\end{pmatrix}\)
\item \(\begin{pmatrix}\frac{1}{2}\\\frac{3}{2}\\3\end{pmatrix}\)
\item \(\begin{pmatrix}1\\3\\3\end{pmatrix}\)
\item \(\begin{pmatrix}\frac{1}{\sqrt{2}}\\\frac{3}{\sqrt{2}}\\\frac{3}{\sqrt{2}}\end{pmatrix}\)
\end{enumerate}
\smallskip
\textbf{3.} \textit{(choose one)} An orthonormal set must be linearly independent.

\begin{enumerate}[label=(\Alph*)]
\item True
\item False
\end{enumerate}
\smallskip
\textbf{4.} \textit{(choose one)} A square matrix \(\mathbf{A}\) of order \(n\) is orthogonal if the columns or rows of \(\mathbf{A}\) form an orthogonal basis for \(\mathbb{R}^n\).

\begin{enumerate}[label=(\Alph*)]
\item False
\item True
\end{enumerate}
\smallskip
\textbf{5.} \textit{(choose one)} Suppose \(\{u_1,u_2,...,u_k\}\) is an orthogonal basis for a subspace \(V\), and \(\mathbf v\) is a vector in \(V\). If \(\mathbf{v}=c_1\mathbf{u}_1+\cdots+c_k\mathbf{u}_k\), then

\begin{enumerate}[label=(\Alph*)]
\item \( c_i=\frac{\mathbf v\cdot \mathbf u_i}{\mathbf u_i\cdot \mathbf u_i}\)
\item \(\ c_i=\frac{\mathbf v\cdot \mathbf u_i}{\mathbf v\cdot \mathbf v}\)
\item \( c_i=(\mathbf v\cdot \mathbf u_i)\)
\item \(\ c_i=\frac{\mathbf v\cdot \mathbf v}{\mathbf u_i\cdot \mathbf u_i}\)
\end{enumerate}
\smallskip
\section{Quiz 5.3}
\textbf{1.} \textit{(choose one)} Which of the following is the set obtained after applying Gram-Schmidt process (with normalization) on the set \(\left\{\begin{pmatrix}1\\1\\0\end{pmatrix},\begin{pmatrix}1\\0\\1\end{pmatrix} \right\}\)?

\begin{enumerate}[label=(\Alph*)]
\item \(\left\{\begin{pmatrix}1/\sqrt{2}\\1/\sqrt{2}\\0\end{pmatrix},\begin{pmatrix}1/\sqrt{6}\\-1/sqrt{6}\\2/\sqrt{6}\end{pmatrix} \right\}\)
\item \(\left\{\begin{pmatrix}1/\sqrt{2}\\1/\sqrt{2}\\0\end{pmatrix},\begin{pmatrix}0\\0\\1\end{pmatrix} \right\}\)
\item \(\left\{\begin{pmatrix}1/\sqrt{2}\\1/\sqrt{2}\\0\end{pmatrix},\begin{pmatrix}1/\sqrt{2}\\-1/\sqrt{2}\\0\end{pmatrix} \right\}\)
\item \(\left\{\begin{pmatrix}1/\sqrt{2}\\1/\sqrt{2}\\0\end{pmatrix},\begin{pmatrix}1/\sqrt{2}\\-1/\sqrt{2}\\1\end{pmatrix} \right\}\)
\end{enumerate}
\smallskip
\textbf{2.} \textit{(choose one)} Which of the following is the set obtained after applying Gram-Schmidt process (with normalization) on the set \(\left\{\begin{pmatrix}1\\1\\2\end{pmatrix},\begin{pmatrix}1\\1\\1\end{pmatrix} \right\}\)?

\begin{enumerate}[label=(\Alph*)]
\item \(\left\{\begin{pmatrix}1/\sqrt{6}\\1/\sqrt{6}\\2/\sqrt{6}\end{pmatrix},\begin{pmatrix}1/\sqrt{3}\\1/\sqrt{3}\\-1/\sqrt{3}\end{pmatrix} \right\}\)
\item \(\left\{\begin{pmatrix}1/\sqrt{6}\\1/\sqrt{6}\\2/\sqrt{6}\end{pmatrix},\begin{pmatrix}1/\sqrt{2}\\-1/\sqrt{2}\\0\end{pmatrix} \right\}\)
\item \(\left\{\begin{pmatrix}1/\sqrt{6}\\1/\sqrt{6}\\2/\sqrt{6}\end{pmatrix},\begin{pmatrix}1/\sqrt{2}\\0\\-1/\sqrt{2}\end{pmatrix} \right\}\)
\item \(\left\{\begin{pmatrix}1/\sqrt{6}\\1/\sqrt{6}\\2/\sqrt{6}\end{pmatrix},\begin{pmatrix}0\\2/\sqrt{5}\\-1/\sqrt{5}\end{pmatrix} \right\}\)
\end{enumerate}
\smallskip
\textbf{3.} \textit{(choose one)} Suppose \(\{u_1,u_2,u_3,u_4\}\) is a linearly independent set and \(\{v_1,v_2,v_3,v_4\}\) is the set obtained after Gram-Schmidt process. Which statement is \textbf{FALSE?}

\begin{enumerate}[label=(\Alph*)]
\item It may be that \(\mathbf{v}_i=\mathbf{0}\) is the zero vector for some i.
\item \(span\{\mathbf{u}_1,\mathbf{u}_2,\mathbf{u}_3\}=span\{\mathbf{v}_1,\mathbf{v}_2,\mathbf{v}_3\}\)
\item \(\mathbf{v}_3\perp span\{\mathbf{u}_1,\mathbf{u}_2\}\)
\item the projection of \(\mathbf{v}_4\) onto \(span\{\mathbf{u}_1,\mathbf{u}_2,\mathbf{u}_3\}\) is the zero vector.
\end{enumerate}
\smallskip
\textbf{4.} \textit{(fill in the blanks)} Suppose a plane \(V\) has a basis \(\left\{\begin{pmatrix}1\\0\\1\end{pmatrix},\begin{pmatrix}0\\1\\0\end{pmatrix}\right\}\). Let \(\mathbf{w}=\begin{pmatrix}4\\2\\6\end{pmatrix}\). Then the orthogonal projection of \(\mathbf{w}\) onto \(V\) is \(\mathbf{w}_p=\begin{pmatrix}x\\y\\z\end{pmatrix}\), where \(x=\) \underline{\hspace{1.2cm}}\textsubscript{1} 
, \(y=\) \underline{\hspace{1.2cm}}\textsubscript{2} 
, \(z=\) \underline{\hspace{1.2cm}}\textsubscript{3} 
.

\smallskip
\textbf{5.} \textit{(choose one)} Let \(S=\{\mathbf{u}_1,\mathbf{u}_2,\mathbf{u}_3,\mathbf{u}_4\}\) be a set in \(\mathbb{R}^4\). Let \(\{\mathbf{v}_1,\mathbf{v}_2,\mathbf{v}_3,\mathbf{v}_4\}\) be the vectors obtained after performing the Gram-Schmidt process on \(S\), before normalization. Suppose \(\mathbf{v}_4=\mathbf{0}\) but \(\mathbf{v}_3\neq\mathbf{0}\). Which of the following statement is true?

\begin{enumerate}[label=(\Alph*)]
\item \(\text{span}(S)\) is a 3-dimensional subspace.
\item The set \(S\) is linearly independent.
\item \(\{\mathbf{v}_1,\mathbf{v}_2,\mathbf{v}_3,\mathbf{v}_4\}\) can be normalized to become an orthonormal set.
\end{enumerate}
\smallskip
\section{Quiz 5.4}
\textbf{1.} \textit{(choose one)} Suppose \(\mathbf{A}\) is a 5 by 4 matrix where the columns are linearly independent. Let \(\mathbf{A=QR}\) be a QR factorization of \(\mathbf{A}\). Which of the following is \textbf{false} about \(\mathbf{R}\)?

\begin{enumerate}[label=(\Alph*)]
\item \(\mathbf{R}\) can have a zero column
\item \(\det(\mathbf{R}) 0\).
\item \(\mathbf{R}\) is a 4 by 4 invertible matrix.
\item \(\mathbf{R}\) is an upper triangular matrix.
\end{enumerate}
\smallskip
\textbf{2.} \textit{(choose all that apply)} Given \(\mathbf{A}=\begin{pmatrix}1&1&0\\1&-1&0\\0&0&1\end{pmatrix}\). Let \(\mathbf{A=QR}\) be a QR factorization of \(\mathbf{A}\). Which of the following statements are true?

\begin{enumerate}[label=(\Alph*)]
\item The \(i\)-th column of \(\mathbf{Q}\) is a multiple of the \(i\)-th column of \(\mathbf{A}\).
\item \(\mathbf{R}\) is a diagonal matrix.
\item None of the other statements are true since \(\mathbf{A}\) is not QR fractorizable.
\item The 3rd column of \(\mathbf{A}\) is equal to the third column of \(\mathbf{Q}\).
\end{enumerate}
\smallskip
\textbf{3.} \textit{(choose all that apply)} Suppose\(\mathbf{A=QR}\) is a QR factorization of \(\mathbf{A}\). Which of the following statement is true?

\begin{enumerate}[label=(\Alph*)]
\item The first column of \(\mathbf{A}\) is a multiple of the first column of \(\mathbf{Q}\).
\item The columns of \(\mathbf{A}\) form an orthogonal set.
\item The columns of \(\mathbf{R}\) form an orthogonal set.
\item None of the other options are true.
\item \(\mathbf{Q}\) is an orthogonal matrix.
\end{enumerate}
\smallskip
\textbf{4.} \textit{(choose one)} Let \(\mathbf{A}=\begin{pmatrix}1&0&0\\0&1&1\\0&0&1\\0&0&0\end{pmatrix}\). which of the following matrices is the matrix \(\mathbf{R}\) in the QR factorization on \(\mathbf{A}\)?

\begin{enumerate}[label=(\Alph*)]
\item \(\begin{pmatrix}1&0&0\\0&1&1\\0&0&1\end{pmatrix}\)
\item \(\begin{pmatrix}1&0&0\\0&1&1\\0&0&-1\end{pmatrix}\)
\item \(\begin{pmatrix}1&0&1\\0&1&0\\0&0&1\end{pmatrix}\)
\item \(\begin{pmatrix}1&0&0\\0&1&0\\0&1&1\end{pmatrix}\)
\item \(\begin{pmatrix}1&0&1\\0&1&1\\0&0&1\end{pmatrix}\)
\end{enumerate}
\smallskip
\textbf{5.} \textit{(choose all that apply)} Suppose \(\mathbf{A=Q}_1\mathbf{R}_1\) is a QR factorization of \(\mathbf{A}\), and  \(2\mathbf{A=Q}_2\mathbf{R}_2\) is a QR factorization of \(2\mathbf{A}\) . Which of the following is true?

\begin{enumerate}[label=(\Alph*)]
\item \(\mathbf{Q}_1=\mathbf{Q}_2\)
\item \(\mathbf{R}_2=2\mathbf{R}_1\)
\item \(\mathbf{Q}_2=2\mathbf{Q}_1\)
\item \(\mathbf{R}_1=\mathbf{R}_2\)
\end{enumerate}
\smallskip
\section{Quiz 5.5}
\textbf{1.} \textit{(choose one)} Suppose \(\{\mathbf u_1,...,\mathbf u_k\}\) is a basis for a subspace \(V\). Let \(\mathbf A=\begin{pmatrix}\mathbf u_1&\cdots&\mathbf u_k\end{pmatrix}\) and \(\mathbf w\) be a vector. Then \(\mathbf w-\mathbf A(\mathbf A^T\mathbf A)^{-1}\mathbf A^T\mathbf w\) is

\begin{enumerate}[label=(\Alph*)]
\item orthogonal to \(V\).
\item the projection of \(\mathbf w\) onto \(V\).
\item neither orthogonal to nor belongs to \(V\).
\end{enumerate}
\smallskip
\textbf{2.} \textit{(choose one)} Let \(\mathbf{u}\) and \(\mathbf{v}\) be solutions to \(\mathbf{A}^T\mathbf{Ax=A}^T\mathbf{b}\). Which of the following statements is true?

\begin{enumerate}[label=(\Alph*)]
\item \(\mathbf{Au=Av}\)
\item \(\mathbf{u}\) and \(\mathbf{v}\) are the projection of \(\mathbf{b}\) onto the column space of\(\mathbf{A}\).
\item \(\mathbf{u=v}\)
\item \(\Vert\mathbf{Au-b}\Vert\neq\Vert\mathbf{Av-b}\Vert\)
\end{enumerate}
\smallskip
\textbf{3.} \textit{(choose one)} If \(\mathbf{u}\) is a solution to \(\mathbf{Ax=b}\), then \(\mathbf{u}\) is a least square solution to \(\mathbf{Ax=b}\).

\begin{enumerate}[label=(\Alph*)]
\item True
\item False
\end{enumerate}
\smallskip
\textbf{4.} \textit{(choose one)} For any \(m\times n\) matrix \(\mathbf{A}\) and \(\mathbf{b}\in\mathbb{R}^m\), \(\mathbf{A}^T\mathbf{Ax}=\mathbf{A}^T\mathbf{b}\) is always consistent.

\begin{enumerate}[label=(\Alph*)]
\item True
\item False
\end{enumerate}
\smallskip
\textbf{5.} \textit{(choose one)} Suppose \(rank(\mathbf{A})=\) number of columns of \(\mathbf{A}\). Which of the statements is always true?

\begin{enumerate}[label=(\Alph*)]
\item \((\mathbf{A}^T\mathbf{A})^{-1}\mathbf{A}^T\mathbf{b}\) is the unique least square solution to \(\mathbf{Ax=b}\).
\item \((\mathbf{A}^T\mathbf{A})^{-1}\mathbf{A}^T\mathbf{b}\) is the unique solution to \(\mathbf{Ax=b}\).
\item \(\mathbf{A}(\mathbf{A}^T\mathbf{A})^{-1}\mathbf{A}^T\mathbf{b}\) is the unique least square solution to \(\mathbf{Ax=b}\).
\item None of the above statement is always true.
\end{enumerate}
\smallskip
\section{Quiz 6.1}
\textbf{1.} \textit{(fill in the blanks)} Suppose the characteristic polynomial of \(\mathbf{A}\) is \(\text{det}(x\mathbf{I}-\mathbf{A})=(x-1)^2(x+3)\). Then the eigenvalue 1 has algebraic multiplicity  \underline{\hspace{1.2cm}}\textsubscript{1} 
 and the eigenvalue -3 has algebraic multiplicity  \underline{\hspace{1.2cm}}\textsubscript{2}

\smallskip
\textbf{2.} \textit{(choose all that apply)} Suppose the characteristic polynomial of \(\mathbf{A}\) is \(\text{det}(x\mathbf{I}-\mathbf{A})=x^2-5x+6\). What are the eigenvalues of \(\mathbf{A}\)?

\begin{enumerate}[label=(\Alph*)]
\item 2
\item 3
\item -3
\item -2
\item -1
\item 1
\end{enumerate}
\smallskip
\textbf{3.} \textit{(choose one)} Let \(\mathbf{A}\) be a 4 by 4 invertible matrix and \(\lambda\) an eigenvalue of \(\mathbf{A}\). which of the following statement is true?

\begin{enumerate}[label=(\Alph*)]
\item There are infinitely many \(\mathbf{v}\) such that \(\mathbf{Av=}\lambda\mathbf{v}\).
\item \(\mathbf{A}\) must have 3 other eigenvalues.
\item \(\lambda\) can be 0
\item The algebraic multiplicity of \(\lambda\) must be 1.
\end{enumerate}
\smallskip
\textbf{4.} \textit{(choose all that apply)} Let \(\mathbf{A}=\begin{pmatrix}5&1&2\\0&6&0\\1&-1&4\end{pmatrix}\). Suppose \(\mathbf{A}\begin{pmatrix}1\\0\\-1\end{pmatrix}=\begin{pmatrix}3\\0\\-3\end{pmatrix}\), \(\mathbf{A}\begin{pmatrix}1\\1\\0\end{pmatrix}=\begin{pmatrix}6\\6\\0\end{pmatrix}\), and \(\mathbf{A}\begin{pmatrix}1\\-1\\1\end{pmatrix}=\begin{pmatrix}6\\-6\\6\end{pmatrix}\). Observe that \(\begin{pmatrix}1\\1\\0\end{pmatrix}+\begin{pmatrix}1\\-1\\1\end{pmatrix}=\begin{pmatrix}2\\0\\1\end{pmatrix}\). Which of the following statements are true?

\begin{enumerate}[label=(\Alph*)]
\item \(\begin{pmatrix}2\\0\\1 \end{pmatrix}\) is an eigenvector associated to 6.
\item \(\begin{pmatrix}2\\0\\1 \end{pmatrix}\) is an eigenvector associated to 3.
\item \(\begin{pmatrix}2\\0\\1 \end{pmatrix}\) is not an eigenvector of \(\mathbf{A}\).
\end{enumerate}
\smallskip
\textbf{5.} \textit{(match)} Let \(\mathbf{A}=\begin{pmatrix}5&1&2\\0&6&0\\1&-1&4\end{pmatrix}\). Suppose \(\mathbf{A}\begin{pmatrix}1\\0\\-1\end{pmatrix}=\begin{pmatrix}3\\0\\-3\end{pmatrix}\), \(\mathbf{A}\begin{pmatrix}1\\1\\0\end{pmatrix}=\begin{pmatrix}6\\6\\0\end{pmatrix}\), and \(\mathbf{A}\begin{pmatrix}1\\-1\\1\end{pmatrix}=\begin{pmatrix}6\\-6\\6\end{pmatrix}\). For each vectors, decide if it is an eigenvector, associated to some eigenvalue, or not.

\begin{enumerate}[label=(\alph*)]
\item \(\begin{pmatrix}1\\0\\-1 \end{pmatrix}\)
\item \(\begin{pmatrix}1\\1\\0 \end{pmatrix}\)
\item \(\begin{pmatrix}1\\-1\\1 \end{pmatrix}\)
\end{enumerate}
\quad Options: 3 \textbar{} 6 \textbar{} not an eigenvector

\smallskip
\section{Quiz 6.2}
\textbf{1.} \textit{(choose one)} Let \(\mathbf{A}\) be a diagonalizable matrix. Then \(\mathbf{A=PDP}^{-1}\) is invertible if and only if

\begin{enumerate}[label=(\Alph*)]
\item the diagonal entries of \(\mathbf{D}\) are nonzero.
\item \(\mathbf{P}\) is invertible.
\item \(\mathbf{D}\) is a nonzero matrix.
\item \(\mathbf{D}\) does not have any zero entries.
\end{enumerate}
\smallskip
\textbf{2.} \textit{(choose all that apply)} Suppose \(\mathbf{A}=\begin{pmatrix}1&1&1\\0&1&1\\0&0&1\end{pmatrix}\begin{pmatrix}1&0&0\\0&2&0\\0&0&2\end{pmatrix}\begin{pmatrix}1&1&1\\0&1&1\\0&0&1\end{pmatrix}^{-1}\). Which of the following product of matrices are \(\mathbf{A}\)?

\begin{enumerate}[label=(\Alph*)]
\item \(\mathbf{A}=\begin{pmatrix}1&1&1\\0&1&1\\0&1&0\end{pmatrix}\begin{pmatrix}1&0&0\\0&2&0\\0&0&2\end{pmatrix}\begin{pmatrix}1&1&1\\0&1&1\\0&1&0\end{pmatrix}^{-1}\)
\item \(\mathbf{A}=\begin{pmatrix}1&1&1\\1&1&0\\1&0&0\end{pmatrix}\begin{pmatrix}2&0&0\\0&2&0\\0&0&1\end{pmatrix}\begin{pmatrix}1&1&1\\1&1&0\\1&0&0\end{pmatrix}^{-1}\)
\item \(\mathbf{A}=\begin{pmatrix}1&1&1\\1&0&1\\1&0&0\end{pmatrix}\begin{pmatrix}2&0&0\\0&1&0\\0&0&2\end{pmatrix}\begin{pmatrix}1&1&1\\1&0&1\\1&0&0\end{pmatrix}^{-1}\)
\item \(\mathbf{A}=\begin{pmatrix}1&1&1\\1&1&0\\1&0&0\end{pmatrix}\begin{pmatrix}1&0&0\\0&2&0\\0&0&2\end{pmatrix}\begin{pmatrix}1&1&1\\1&1&0\\1&0&0\end{pmatrix}^{-1}\)
\end{enumerate}
\smallskip
\textbf{3.} \textit{(choose one)} Let \(\lambda\) be an eigenvalue associated to square matrix \(\mathbf A\), \(r_\lambda\) its algebraic multiplicity, and \(E_\lambda\) its eigenspace. Which of the following inequalities is true?

\begin{enumerate}[label=(\Alph*)]
\item \(dim(E_\lambda)\leq r_\lambda\).
\item \(r_\lambda\leq \dim(E_\lambda)\).
\item There is no relations between \(dim(E_\lambda)\) and \(r_\lambda\).
\end{enumerate}
\smallskip
\textbf{4.} \textit{(choose one)} Which of the following matrix \(\mathbf{A}\) is \textbf{not} diagonalizable.

\begin{enumerate}[label=(\Alph*)]
\item \(\text{det}(\lambda\mathbf{I}-\mathbf{A})=(\lambda-2)^2(\lambda-1)\) and the eigenspace \(E_2\) is 1-dimensional.
\item \(\mathbf{A}\) is an order 3 square matrix with 2 eigenvalues \(\lambda_1\) and \(\lambda_2\) such that the eigenspace \(E_{\lambda_1}\) is 2-dimensional.
\item \(\mathbf{A}\begin{pmatrix}1\\1\\1\end{pmatrix}=2\begin{pmatrix}1\\1\\1\end{pmatrix}\), \(\mathbf{A}\begin{pmatrix}1\\1\\0\end{pmatrix}=\begin{pmatrix}1\\1\\0\end{pmatrix}\), \(\mathbf{A}\begin{pmatrix}1\\0\\1\end{pmatrix}=\begin{pmatrix}1\\0\\1\end{pmatrix}\).
\item \(\mathbf{A}\) is an order 3 square matrix with 3 distinct eigenvalues,
\end{enumerate}
\smallskip
\textbf{5.} \textit{(choose one)} Let \(\mathbf{A}=\begin{pmatrix}3&-1&-1\\1&1&-1\\1&-1&1 \end{pmatrix}\), \(\mathbf{P}=\begin{pmatrix}1&1&1\\1&1&0\\1&0&1 \end{pmatrix}\), and \(\mathbf{D}=\begin{pmatrix}1&0&0\\0&2&0\\0&0&2\end{pmatrix}\). Given that \(\mathbf{A}=\mathbf{PDP}^{-1}\). Which statement is true?

\begin{enumerate}[label=(\Alph*)]
\item \(\begin{pmatrix}1\\1\\0\end{pmatrix}\) is an eigenvector of \(\mathbf{A}\) associated to eigenvalue 2.
\item \(\begin{pmatrix}1\\0\\1\end{pmatrix}\) is an eigenvector of \(\mathbf{A}\) associated to eigenvalue 1.
\item The eigenspace \(E_2\) assocaited to eigenvalue 2 is 1-dimensional.
\end{enumerate}
\smallskip
\section{Quiz 6.3}
\textbf{1.} \textit{(choose all that apply)} Suppose \(\mathbf{A}\) is a 3 by 3 symmetric matrix such that \(\mathbf{P}^T\mathbf{AP=D}\), where \(\mathbf{D}\) is a diagonal matrix. Write \(P=\begin{pmatrix}\mathbf{p}_1&\mathbf{p}_2&\mathbf{p}_3\end{pmatrix}\)

Which of the following are equal to \(\mathbf{A}\)?

\begin{enumerate}[label=(\Alph*)]
\item \(\begin{pmatrix}-\mathbf{p}_1\mathbf{p}_2-\mathbf{p}_3\end{pmatrix}\mathbf{D}\begin{pmatrix}-\mathbf{p}_1\mathbf{p}_2-\mathbf{p}_3\end{pmatrix}^T\)
\item \(\begin{pmatrix}1\mathbf{p}_12\mathbf{p}_23\mathbf{p}_3\end{pmatrix}\mathbf{D}\begin{pmatrix}1\mathbf{p}_12\mathbf{p}_23\mathbf{p}_3\end{pmatrix}^{-1}\)
\item \(\begin{pmatrix}1\mathbf{p}_12\mathbf{p}_23\mathbf{p}_3\end{pmatrix}\mathbf{D}\begin{pmatrix}1\mathbf{p}_12\mathbf{p}_23\mathbf{p}_3\end{pmatrix}^{T}\)
\item \(\begin{pmatrix}-\mathbf{p}_1\mathbf{p}_2-\mathbf{p}_3\end{pmatrix}\mathbf{D}\begin{pmatrix}-\mathbf{p}_1\mathbf{p}_2-\mathbf{p}_3\end{pmatrix}^{-1}\)
\end{enumerate}
\smallskip
\textbf{2.} \textit{(choose one)} Which of the following statements is \textbf{FALSE}?

\begin{enumerate}[label=(\Alph*)]
\item If A is 3 by 3 diagonalizable matrix, then A must have 3 different eigenvalues.
\item A 3 by 3 matrix with eigenvalues -1, 0, 2, is diagonalizable.
\item A 5 by 5 symmetric matrix is diagonalizable.
\item If A is symmetric, then all the eigenspaces are orthogonal to each other.
\end{enumerate}
\smallskip
\textbf{3.} \textit{(choose one)} Suppose a 3 by 3 symmetric matrix \(\mathbf{A}\) has \(\begin{pmatrix}1\\1\\1\end{pmatrix}\), \(\begin{pmatrix}1\\0\\0\end{pmatrix}\), and \(\begin{pmatrix}0\\1\\-1\end{pmatrix}\) as eigenvectors (not necessarily associated to the same eigenvalue). Which of the statements is false?

\begin{enumerate}[label=(\Alph*)]
\item \(\mathbf{A}\) can have 3 distinct eigenvalues.
\item \(\mathbf{A}\) can have 2 distinct eigenvalues.
\item \(\mathbf{A}\) can have 1 distinct eigenvalue.
\end{enumerate}
\smallskip
\textbf{4.} \textit{(choose one)} Let \(\mathbf{A}\) be a diagonalizable matrix. By applying Gram-Schmidts process to the bases of the eigenspaces then normalizing them,  we can always orthogonally diagonalize \(\mathbf{A}\).

\begin{enumerate}[label=(\Alph*)]
\item False
\item True
\end{enumerate}
\smallskip
\textbf{5.} \textit{(choose one)} Suppose a 3 by 3 matrix \(\mathbf{A}\) has eigenvectors \(\begin{pmatrix}1\\1\\1\end{pmatrix}\), \(\begin{pmatrix}0\\2\\2\end{pmatrix}\), and \(\begin{pmatrix}0\\0\\3\end{pmatrix}\). Then \(\mathbf{A}\) cannot be a symmetric matrix.

\begin{enumerate}[label=(\Alph*)]
\item False
\item True
\end{enumerate}
\smallskip
\section{Quiz 6.4}
\textbf{1.} \textit{(choose one)} In a population of rabbits,

\quad\textbullet\ half of the newborn rabbits survive their first year;

\quad\textbullet\ of those, half survive their second year;

\quad\textbullet\ the maximum life span is three years;

\quad\textbullet\ rabbits produce  0 ,  6 ,  8  rabbits in their first, second, and third years, respectively.

Let \(\mathbf{v}_n=\begin{pmatrix}x_n\\y_n\\z_n\end{pmatrix}\), where \(x_n,y_n,z_n\) are the number of rabbits aged  0, 1, and  2, respectively. Which of the following is the transition matrix?

\begin{enumerate}[label=(\Alph*)]
\item \(\begin{pmatrix}068\\1/200\\01/20\end{pmatrix}\)
\item \(\begin{pmatrix}01/20\\601/2\\800\end{pmatrix}\)
\item \(\begin{pmatrix}068\\01/20\\1/200\end{pmatrix}\)
\item \(\begin{pmatrix}001/2\\61/20\\800\end{pmatrix}\)
\end{enumerate}
\smallskip
\textbf{2.} \textit{(choose all that apply)} Which of the following vectors are probability vectors?

\begin{enumerate}[label=(\Alph*)]
\item \(\begin{pmatrix}0\\0\\1\end{pmatrix}\)
\item \(\begin{pmatrix}0.5\\0.2\\0.3\end{pmatrix}\)
\item \(\begin{pmatrix}0.5\\0.6\\-0.1\end{pmatrix}\)
\item \(\begin{pmatrix}0.5\\0.6\\0.2\end{pmatrix}\)
\end{enumerate}
\smallskip
\textbf{3.} \textit{(choose one)} Suppose there are 3 cities, A, B, and C in a certain country. It is known that no one migrates in and out of the country, and that the birth and death rates are equal. Every year, 5\textbackslash{}\% of the population in city A will move to city B and 3\textbackslash{}\% will move to city C, 2\textbackslash{}\% of the population in city B will move to city A and 2\textbackslash{}\% will move to city C, and 7\textbackslash{}\% of the population of city C will move to city A and 1\textbackslash{}\% of the population will move to city B. In the long run, what will be the approximate population distribution?

\begin{enumerate}[label=(\Alph*)]
\item 32\textbackslash{}\% of the population of the country will live in city A, 45\textbackslash{}\% will live in city B, and 23\textbackslash{}\% will live in city C.
\item 48\textbackslash{}\% of the population of the country will live in city A, 31\textbackslash{}\% will live in city B, and 21\textbackslash{}\% will live in city C.
\item 22\textbackslash{}\% of the population of the country will live in city A, 44\textbackslash{}\% will live in city B, and 34\textbackslash{}\% will live in city C.
\item There are not enough information to determine the population distribution.
\end{enumerate}
\smallskip
\textbf{4.} \textit{(choose one)} Which of the follow matrices is stochastic?

\begin{enumerate}[label=(\Alph*)]
\item \(\begin{pmatrix}0.2&0.6&0.3\\0.3&0.3&0.3\\0.5&0.1&0.4\end{pmatrix}\)
\item \(\begin{pmatrix}0.2&0.3&0.5\\0.6&0.3&0.1\\0.3&0.3&0.4\end{pmatrix}\)
\item \(\begin{pmatrix}0.2&0.3&-0.5\\0.3&0.5&1.3\\0.5&0.2&0.2\end{pmatrix}\)
\end{enumerate}
\smallskip
\textbf{5.} \textit{(choose one)} Every year a certain company will hire one of the three catering companies, caterer A, caterer B, or caterer C, to prepare the food for their company retreat. The company never hire the same caterer twice in a row. If the company hires caterer A this year, there is a \(45\%\) chance it will hire caterer B the next year. If the company hires caterer B this year, there is a \(30\%\) chance it will hire caterer C the next year. If the company hires caterer C this year, there is a \(50\%\) chance it will hire caterer A the next year. What is the probability transition matrix?

\begin{enumerate}[label=(\Alph*)]
\item \(\begin{pmatrix}0&0.7&0.5\\0.45&0&0.5\\0.55&0.3&0\end{pmatrix}\)
\item \(\begin{pmatrix}0&0.45&0.55\\0.7&0&0.3\\0.5&0.5&0\end{pmatrix}\)
\item \(\begin{pmatrix}0&0.45&0.7\\0.5&0&0.3\\0.5&0.55&0\end{pmatrix}\)
\item There are not enough information given to find the probability transition matrix.
\end{enumerate}
\smallskip
\section{Quiz 6.5}
\textbf{1.} \textit{(choose for each blank)} Let \(\mathbf{A=U\Sigma V}^T\) be a SVD decomposition of a \(m\times n\) matrix \(\mathbf{A}\). Then \(rank(\mathbf{U})=\) \underline{\hspace{1.2cm}}\textsubscript{1} 
, \(rank(\mathbf{\Sigma})=\) \underline{\hspace{1.2cm}}\textsubscript{2} 
, and \(rank(\mathbf{V})=\) \underline{\hspace{1.2cm}}\textsubscript{3} 
.

\quad Blank 1: m \textbar{} n \textbar{} \(rank(\mathbf{A})\) \textbar{} None of the other options

\quad Blank 2: \(rank(\mathbf{A})\) \textbar{} n \textbar{} m \textbar{} None of the other options

\quad Blank 3: n \textbar{} m \textbar{} \(rank(\mathbf{A})\) \textbar{} None of the other options

\smallskip
\textbf{2.} \textit{(choose all that apply)} Let \(\mathbf{A}\) be a \(m\times n \) matrix and \(\mathbf{U\Sigma V}^T\) be a singular value decomposition of \(\mathbf{A}\). Which of the following statements are true?

\begin{enumerate}[label=(\Alph*)]
\item \(\mathbf{U}\) is a order \(m\) square matrix
\item \(\mathbf{V}\) is a order \(n\) square matrix
\item \(\mathbf{\Sigma}\) is a diagonal matrix.
\item \(\mathbf{U}\) has the same size as \(\mathbf{A}\)
\end{enumerate}
\smallskip
\textbf{3.} \textit{(choose one)} Let \(\mathbf{A=U\Sigma V}^T\) be a SVD decomposition of \(\mathbf{A}\). Define \(\mathbf{\Sigma}_t\) to be the matrix obtained from \(\mathbf{\Sigma}\) by replacing the \(i\)-th singular value \(\sigma_i\) with \(0\) if \(i\geq t\). 

Define \(\mathbf{A}_t=\mathbf{U\Sigma}_t\mathbf{V}^T\). Then it must be that

\begin{enumerate}[label=(\Alph*)]
\item \(rank(\mathbf{A})\leq rank(\mathbf{A}_t)\).
\item \(rank(\mathbf{A}) rank(\mathbf{A}_t)\).
\item \(rank(\mathbf{A})=rank(\mathbf{A}_t)\).
\item \(rank(\mathbf{A})\geq rank(\mathbf{A}_t)\).
\item \(rank(\mathbf{A})rank(\mathbf{A}_t)\).
\end{enumerate}
\smallskip
\textbf{4.} \textit{(choose one)} SVD can be applied to

\begin{enumerate}[label=(\Alph*)]
\item any matrix.
\item only full rank matrices.
\item only square matrices.
\item only diagonalizable matrices.
\item only invertible matrices.
\item only orthogonal matrices.
\end{enumerate}
\smallskip
\textbf{5.} \textit{(choose for each blank)} Let \(\mathbf{A}=\begin{pmatrix}1&0&2\\0&1&2\end{pmatrix}\). It is given that \(det(x\mathbf{I-A}^T\mathbf{A})=x(x-1)(x-9)\). Let\(\mathbf{A=U\Sigma V}^T\) be a singular value decomposition of \(\mathbf{A}\). Then the matrix \(\mathbf{\Sigma}\) is a  \underline{\hspace{1.2cm}}\textsubscript{1} 
 matrix, the (1,1)-entry is  \underline{\hspace{1.2cm}}\textsubscript{2} 
 and the (2,2)-entry is  \underline{\hspace{1.2cm}}\textsubscript{3}

\quad Blank 1: 2 by 3 \textbar{} 3 by 2 \textbar{} 2 by 2 \textbar{} 3 by 3

\quad Blank 2: 3 \textbar{} 0 \textbar{} 1 \textbar{} 9

\quad Blank 3: 1 \textbar{} 0 \textbar{} 3 \textbar{} 9

\smallskip
\section{Quiz 7.1}
\textbf{1.} \textit{(choose for each blank)} Let \(T:\mathbb{R}^m\to\mathbb{R}^n\) be a matrix transformation, \(T(\mathbf{u})=\mathbf{Au}\) for all \(\mathbf{u}\in\mathbb{R}^m\). Then \(\mathbf{A}\) is a  \underline{\hspace{1.2cm}}\textsubscript{1} 
 by  \underline{\hspace{1.2cm}}\textsubscript{2} 
 matrix.

\quad Blank 1: n \textbar{} m

\quad Blank 2: m \textbar{} n

\smallskip
\textbf{2.} \textit{(choose one)} A mapping \(T;\mathbb{R}^3\to\mathbb{R}^1\) is such that \(T\left(\begin{pmatrix}0\\0\\0\end{pmatrix} \right)=0\), \(T\left(\begin{pmatrix}1\\0\\0\end{pmatrix} \right)=5\), and \(T\left(\begin{pmatrix}2\\0\\0\end{pmatrix} \right)=6\). What can we conclude?

\begin{enumerate}[label=(\Alph*)]
\item None of the other options
\item There is a \(3\times 1\) matrix \(\mathbf{A}\) such that \(T(\mathbf{u})=\mathbf{Au}\) for all \(\mathbf{u}\) in \(\mathbb{R}^3\).
\item There is a \(1\times 3\) matrix \(\mathbf{A}\) such that \(T(\mathbf{u})=\mathbf{Au}\) for all \(\mathbf{u}\) in \(\mathbb{R}^3\).
\end{enumerate}
\smallskip
\textbf{3.} \textit{(choose all that apply)} Which of the following mapping is a linear transformation?

\begin{enumerate}[label=(\Alph*)]
\item \(T\left(\begin{pmatrix}x\\y\\z\end{pmatrix}\right)=\begin{pmatrix}y\\x\end{pmatrix}\)
\item \(T\left(\begin{pmatrix}x\\y\end{pmatrix}\right)=\begin{pmatrix}0\\x\\y\end{pmatrix}\)
\item \(T\left(\begin{pmatrix}x\\y\end{pmatrix}\right)=\begin{pmatrix}x^2\\y^2\\0\end{pmatrix}\)
\item \(T\left(\begin{pmatrix}x\\y\\z\end{pmatrix}\right)=x\begin{pmatrix}1\\1\\z\end{pmatrix}\)
\item \(T\left(\begin{pmatrix}x\\y\end{pmatrix}\right)=\begin{pmatrix}1\\x\\y\end{pmatrix}\)
\end{enumerate}
\smallskip
\textbf{4.} \textit{(choose one)} A transformation \(T:\mathbb{R}^3\to\mathbb{R}\) is given by \(T\left(\begin{pmatrix}r\\g\\b\end{pmatrix} \right)=0.299r+0.587g+0.114b\). What is the standard matrix of \(T\)?

\begin{enumerate}[label=(\Alph*)]
\item \(\begin{pmatrix}0.2990.5870.114\end{pmatrix}\)
\item \(\begin{pmatrix}0.299\\0.587\\0.114\end{pmatrix}\)
\item \(\begin{pmatrix}0.1140.5870.299\end{pmatrix}\)
\item \(\begin{pmatrix}0.114\\0.587\\0.299\end{pmatrix}\)
\end{enumerate}
\smallskip
\textbf{5.} \textit{(choose one)} Let \(T\left(\begin{pmatrix}x\\y\end{pmatrix}\right)=\begin{pmatrix}x+y\\x-y\\2y\end{pmatrix}\). The standard matrix of \(T\) is

\begin{enumerate}[label=(\Alph*)]
\item \(\begin{pmatrix}1&1\\1&-1\\0&2\end{pmatrix}\)
\item \(\begin{pmatrix}1&1&0\\1&-1&2\end{pmatrix}\)
\item \(\begin{pmatrix}1&1\\-1&1\\2&0\end{pmatrix}\)
\item \(\begin{pmatrix}1&-1&2\\1&1&0\end{pmatrix}\)
\end{enumerate}
\smallskip
\section{Quiz 7.2}
\textbf{1.} \textit{(choose one)} Suppose the linear transformation \(T:\mathbb{R}^n\to\mathbb{R}^m\) is one-to-one. What can we conclude?

\begin{enumerate}[label=(\Alph*)]
\item The rank of \(T\) is \(n\).
\item The rank of \(T\) is \(m\).
\item The rank of \(T\) is strictly smaller than \(m\) and \(n\).
\item We cannot conclude anything.
\end{enumerate}
\smallskip
\textbf{2.} \textit{(choose one)} Suppose \(T:\mathbb{R}^3\to\mathbb{R}^5\) is a linear transformation such that the standard matrix is full rank. What can we conclude?

\begin{enumerate}[label=(\Alph*)]
\item \(T\) is injective
\item \(T\) is surjective
\item \(T\) is bijective
\item We cannot make any conclusion.
\end{enumerate}
\smallskip
\textbf{3.} \textit{(choose one)} Suppose the linear transformation \(T:\mathbb{R}^n\to\mathbb{R}^m\) is onto. What can we conclude?

\begin{enumerate}[label=(\Alph*)]
\item The rank of \(T\) is \(m\).
\item The rank of \(T\) is \(n\).
\item The rank of \(T\) is strictly smaller than \(m\) and \(n\).
\item We cannot conclude anything.
\end{enumerate}
\smallskip
\textbf{4.} \textit{(choose one)} Suppose \(T:\mathbb{R}^5\to\mathbb{R}^3\) is a linear transformation such that the standard matrix is full rank. What can we conclude?

\begin{enumerate}[label=(\Alph*)]
\item \(T\) is surjective
\item \(T\) is injective
\item \(T\) is bijective
\item We cannot make any conclusion.
\end{enumerate}
\smallskip
\textbf{5.} \textit{(choose one)} Suppose \(T:\mathbb{R}^n\to\mathbb{R}^m\) is a linear transformation such that \(rank(T)=n\). What can we conclude?

\begin{enumerate}[label=(\Alph*)]
\item \(T\) must be one-to-one.
\item \(T\) must be onto.
\item \(T\) is neither one-to-one nor onto.
\item We cannot conclude anything.
\end{enumerate}
\smallskip
\end{document}